% Chapitre 16 — L'analyse inter-composants et inter-fichiers
% Sources : notes/08 (intégralement), notes/05 (§4.1.10, §4.2), notes/06 (§4.6).
% Commit photographié : e67b10a.
% Toutes les sorties CLI ont été rejouées sous /tmp/ch16/ avec target/debug/reactant
% construit depuis e67b10a. Les résumés RenderDeps champ par champ (genuine, sites,
% handlers) viennent de la sonde de débogage du dossier 08 (§6), qui n'a pas pu être
% reconstruite pour ce chapitre (/tmp plein) : ils sont signalés comme tels dans le texte.

\chapter{L'analyse inter-composants et inter-fichiers}
\label{chap:inter-composants}

\begin{objectifs}
  \item Comprendre pourquoi un analyseur \emph{intra}-composant, qui lit les props comme
        $\Top$, manque les défauts qui traversent une frontière de composant, et ce qu'il faut
        pour les voir.
  \item Connaître l'identité d'un composant : la clé \code{(file, name)}, le
        \code{ComponentId} interné, le nom d'affichage qui n'est qu'un rendu, l'origine
        \code{CompOrigin} qu'un fichier prouve pour un élément JSX, et les trois réponses de
        \code{resolve_child} (\adr{40}).
  \item Savoir comment les racines sont choisies (\code{RootStrategy}, \code{--all-roots},
        \code{--entry}) et pourquoi une référence non résolue compte pour tous les homonymes.
  \item Suivre pas à pas l'inlining descendant (\code{eval_comp_app}, \code{InterCtx}, pile de
        récursion, props dans le tas) et le flux ascendant (setter porteur de propriétaire,
        store partagé monotone, havoc).
  \item Comprendre le cache de composants, pourquoi il n'est qu'un \og déjà vu\fg{}, et le faux
        négatif qui en découle ; les deux phases d'\code{analyze_program} et la discipline
        \code{phase1_reached}.
  \item Connaître ce que lit une règle (\code{AnalysisResult}, \code{ProgramAnalysisResult},
        \code{RuleCtx}) et ce que le graphe de symboles d'\adr{13} est devenu.
  \item Maîtriser la \terme{dépendance de rendu} (\adr{41}) : le treillis \code{Deps}, la
        dépendance de contrôle, la séparation \code{genuine} / \code{sites}, handlers,
        contextes, bindings de module, fait \code{gated}, et sa composition sur l'arbre
        d'éléments (\code{uses}, \code{home_of}, \code{wasted_siblings}).
\end{objectifs}

\prerequis{\cref{chap:point-fixe-composant} (la boucle externe d'un composant, sa passe de
rendu, le test de convergence, la passe de rafraîchissement) ; \cref{chap:transfert-stores}
(\code{eval_comp_app} vu du transfert, le \code{SharedStateStore}, le tas, le havoc des
setters, \cref{def:store}) ; \cref{chap:statevalue-stabilite} (\cref{def:statevalue}, la
composante setter du produit, \code{InterCtx}) ; \cref{chap:ir} (\cref{def:slot},
\cref{def:site}, \cref{def:cfg}, \code{Expr::CompApp}) ; \cref{chap:interpretation-abstraite}
(\cref{def:treillis}, \cref{def:must-may}, \cref{def:soundness}). Les règles qui consomment la
dépendance de rendu sont traitées aux \cref{chap:regles-etat,chap:regles-rendu-rsc}.}

Jusqu'ici, le moteur a été décrit composant par composant : un CFG de rendu, des effets, des
handlers, une boucle externe qui monte l'état jusqu'à un point fixe. Or un programme React est
un \emph{arbre} de composants, écrit dans des dizaines de fichiers, où l'état d'un parent
descend en props et où ses setters descendent pour être appelés plus bas. Ce chapitre répond à
deux questions qu'aucune analyse d'un composant isolé ne peut trancher.

\begin{enumerate}
  \item \textbf{Qui rend qui, avec quelles props ?} C'est l'analyse \emph{inter-composants}
        (\adr{12}) : un parent est analysé d'abord, et quand son rendu rencontre
        \code{<Child a={x}/>}, l'enfant est analysé avec les props \emph{abstraites}
        évaluées dans l'environnement du parent. Elle demande de savoir quel composant un nom
        JSX désigne, à travers les fichiers (\adr{13}, \adr{40}), par où commencer, et comment
        ne pas ré-analyser un enfant pour des props identiques.
  \item \textbf{Quelles parties de la sortie d'un composant peuvent être calculées à partir
        d'un état donné ?} C'est la \emph{dépendance de rendu} (\adr{41}), une seconde
        analyse, séparée, lancée après convergence, qui alimente les règles de cascades de
        re-rendus.
\end{enumerate}

Le sous-système tient dans sept fichiers de \file{src/engine/} :
\file{component_registry.rs} (256 lignes), \file{root_detector.rs} (329),
\file{component_cache.rs} (327), \file{symbol_graph.rs} (416), \file{analysis_result.rs} (289),
\file{program_result.rs} (225) et \file{render_deps.rs} (1\,276). Le mécanisme lui-même est porté
par quatre fichiers voisins : \file{src/ir/component_id.rs} (l'identité),
\file{src/domains/context.rs} (\code{InterCtx}), \file{src/domains/transfer/state_value.rs}
(\code{eval_comp_app}) et \file{src/engine/fixpoint.rs} (\code{analyze_program}). La
composition des résumés de rendu sur l'arbre d'éléments vit dans la couche règles,
\file{src/rules/helpers/render_tree.rs} (777 lignes). Toutes les sorties de l'analyseur
montrées ici ont été observées avec le binaire du commit \snapshotCommit{} sur des fichiers
de \file{/tmp/ch16/} ou sur les fixtures du dépôt.

% ===========================================================================
\section{Le problème : des props inconnues, un setter qui remonte}
\label{sec:inter-composants-probleme}

\subsection{Un badge qui relance son parent}

\begin{exemple}{Un setter appelé pendant le rendu d'un enfant}{inter-composants-badge}
\begin{lstlisting}[language=TypeScript,caption={\file{/tmp/ch16/intro.tsx}},label={lst:inter-composants-intro}]
import { useState } from "react";

function Badge({ onSeen }: { onSeen: (n: number) => void }) {
  onSeen(1);
  return <span>new</span>;
}

export function Inbox() {
  const [seen, setSeen] = useState(0);
  return <div>{seen}<Badge onSeen={setSeen} /></div>;
}
\end{lstlisting}
\end{exemple}

\code{Badge} appelle, pendant son rendu, la fonction que son parent lui passe. Or cette
fonction est \code{setSeen}, le setter d'un slot d'\code{Inbox}. Chaque rendu de \code{Badge}
écrit donc l'état de son parent pendant un rendu, ce que React signale en développement
(\og Cannot update a component while rendering a different component\fg) ; chaque écriture qui
change la valeur re-rend le parent, donc l'enfant, qui écrit de nouveau. L'analyseur le voit :

\begin{sortie}
$ reactant check intro.tsx --trace
  Badge  (0 hooks)  intro.tsx
    error  cross-setter-in-render  var:onSeen  (line 4:2)  prop `onSeen` (a state setter of parent `Inbox`) called during render of `Badge`, which triggers a parent re-render on every render
       → `onSeen` is a state setter, so calling it writes state (line 4:2)
   1 clean component(s) hidden, rerun with --show-clean

⚠  1 error(s) across 1 file(s).
\end{sortie}

Le diagnostic est une \Error{} : l'appel est inconditionnel dans le corps du rendu, et le
message nomme le propriétaire du setter, \code{Inbox}. Ce nom n'est écrit nulle part dans
\code{Badge}.

\subsection{Ce qu'un analyseur intra ne voit pas}

Retirons le parent. Le fichier \file{/tmp/ch16/badge-alone.tsx} ne contient que la fonction
\code{Badge}, verbatim :

\begin{sortie}
$ reactant check badge-alone.tsx --info

✓  1 file(s) no issues found.
\end{sortie}

Aucune alarme, et à juste titre. Analysé seul, \code{Badge} reçoit un objet de props dont
l'analyse ne sait rien : son paramètre n'est lié à rien dans l'environnement d'entrée, et une
variable absente se lit $\Top$ (\code{AbstractEnv::lookup},
\src{src/domains/stores/abstract_env.rs}{78}{81}). \code{onSeen} est donc \og une valeur
quelconque\fg{} ; rien ne prouve que c'est un setter, et une règle \Error{} doit prouver
\emph{toute} sa conclusion. Un analyseur intra-composant est sound (il ne conclut rien de faux),
mais il est aveugle à toute une famille de défauts : un setter passé en prop et appelé
pendant le rendu, un effet d'enfant qui relance le parent (\regle{cross-component-infinite-loop}),
un état qu'un parent possède pour le seul usage d'un arrière-petit-enfant.

\begin{react}[props, setters, identité]
Un setter renvoyé par \code{useState} a une identité stable pour toute la vie du composant,
et il écrit toujours le slot de son \emph{propriétaire}, quel que soit le composant qui
l'appelle. Le passer en prop est l'idiome normal pour qu'un enfant remonte une valeur. Un
\emph{élément} \code{<Badge onSeen={setSeen}/>} n'est qu'une description (\code{jsx(Badge,
props)}) : \code{Badge} ne s'exécute que lorsque React monte l'élément, avec ces props
(\cref{chap:react-modele}).
\end{react}

Pour conclure, l'analyse doit donc savoir que \code{<Badge/>} dans \code{Inbox} désigne
\emph{ce} \code{Badge}-là, analyser \code{Badge} avec la valeur abstraite que prend
\code{onSeen} dans \code{Inbox} --- un setter porteur de son propriétaire --- puis faire
remonter l'écriture de \code{Badge} dans l'état d'\code{Inbox}. C'est exactement ce que fait
l'analyse inter-composants.

\subsection{Deux questions, un sous-système}

La seconde question naît d'un autre programme, la fixture \file{drill.tsx} du répertoire
\file{tests/fixtures/state_lifted_too_high/} :

\tsxsnippet[label={lst:inter-composants-drill}]{tests/fixtures/state_lifted_too_high/drill.tsx}{6}{19}{un état tenu trois niveaux trop haut}

Rien n'y est faux : c'est un problème de coût. À chaque frappe, \code{setValue} re-rend
\code{App}, et donc tous les éléments de composant qu'\code{App} construit, récursivement :
\code{Layout}, \code{Sidebar}, \code{Field}, et \code{Content} au passage. Seul \code{Field}
utilise la valeur. Pour le dire, l'analyse doit savoir, pour chaque composant, ce qui dans sa
sortie \emph{dépend} de l'état et ce qui ne fait que le \emph{transmettre} à un enfant. La
valeur abstraite ne le dit pas (nous verrons pourquoi au \cref{sec:inter-composants-render-deps}) :
il faut une analyse de dépendance.

\begin{figure}[htbp]\centering
\begin{tikzpicture}
  \node[bloc] (files) at (0,0) {fichiers\\\code{.ts}/\code{.tsx}};
  \node[bloc] (low) at (3.2,0) {\code{lower_files_with}\\{\scriptsize IR, \code{CompApp.origin}}};
  \node[bloc] (reg) at (7.0,0) {\code{ComponentRegistry}\\{\scriptsize \code{ComponentTable}}};
  \node[bloc] (roots) at (10.8,0) {\code{RootStrategy}\\\code{::detect}};
  \node[bloc] (prog) at (10.8,-1.9) {\code{analyze_program}\\{\scriptsize phase 1 : racines sous \code{InterCtx}}\\{\scriptsize phase 2 : balayage intra}};
  \node[bloc] (eca) at (10.8,-4.0) {\code{eval_comp_app}\\{\scriptsize inlining d'un enfant}};
  \node[bloc] (res) at (1.2,-1.9) {\code{ProgramAnalysisResult}};
  \node[bloc] (cache) at (1.2,-4.0) {\code{ProgramCache}\\{\scriptsize relations, \code{RenderIndex}}};
  \node[bloc] (rd) at (5.0,-4.0) {\code{render_deps}\\{\scriptsize un résumé par composant}};
  \node[bloc] (rules) at (8.2,-4.0) {\code{RuleCtx}\\\code{Rule::check}};
  \draw[fleche] (files) -- (low);
  \draw[fleche] (low) -- (reg);
  \draw[fleche] (reg) -- (roots);
  \draw[fleche] (roots) -- (prog);
  \draw[flechemust] ($(prog.south)+(-0.4,0)$) -- node[etiquette,left] {\code{<Child/>}} ($(eca.north)+(-0.4,0)$);
  \draw[flechemust] ($(eca.north)+(0.4,0)$) -- node[etiquette,right] {enfant analysé} ($(prog.south)+(0.4,0)$);
  \draw[fleche] (prog) -- (res);
  \draw[fleche] (res) -- (cache);
  \draw[fleche] (cache) -- (rd);
  \draw[fleche] (rd) -- (rules);
\end{tikzpicture}
\caption{Le pipeline programme. L'IR de tous les fichiers est abaissée d'abord ; le registre
interne l'identité des composants ; la stratégie choisit les racines ; \code{analyze_program}
analyse les racines en inlinant leurs enfants à la volée (\code{eval_comp_app}), puis balaie
les composants non atteints. Les résumés de dépendance de rendu sont calculés à la demande,
une fois par programme, dans la couche règles.}
\label{fig:inter-composants-pipeline}
\end{figure}

La \cref{fig:inter-composants-pipeline} situe les deux mécanismes. Les points d'entrée exacts
sont \code{resolver::analyze_lowered} (\src{src/resolver/mod.rs}{439}{452}), qui construit les
registres et appelle \code{engine::analyze_program}
(\src{src/engine/fixpoint.rs}{704}{819}) ; puis le driver construit un \code{ProgramCache}
partagé par tous les \code{RuleCtx}, dont la première demande de \code{render()} calcule les
résumés \code{render_deps} de tous les composants. Retenons un point dès maintenant : dans le
domaine des valeurs, un élément JSX de composant vaut \emph{toujours}
\code{StateValue::reference(Stability::Stable)} ; ce qui compte, ce sont les effets de bord de
son évaluation --- le résultat de l'enfant, les écritures dans le store partagé, les arêtes du
graphe d'appel.

% ===========================================================================
\section{L'identité d'un composant}
\label{sec:inter-composants-identite}

\subsection{Trois façons de nommer un composant}

Avant de pouvoir analyser un enfant, il faut savoir \emph{lequel}. Considérons quatre
fichiers (\file{/tmp/ch16/amb/}) :

\begin{lstlisting}[language=TypeScript,caption={Deux \code{Widget}, un alias, une référence nue (\file{/tmp/ch16/amb/})},label={lst:inter-composants-amb}]
// a/Widget.tsx
import { useEffect, useState } from "react";
export function Widget() {
  const [n, setN] = useState(0);
  useEffect(() => { setN(n + 1); });
  return <p>{n}</p>;
}
// b/Widget.tsx
export function Widget({ onReady }: { onReady: (v: number) => void }) {
  onReady(1);
  return <p>b</p>;
}
// App.tsx
import { useState } from "react";
import { Widget as Panel } from "./b/Widget";
export function App() {
  const [v, setV] = useState(0);
  return <Panel onReady={setV} />;
}
// Page.tsx  (aucun import de Widget)
import { useState } from "react";
export function Page() {
  const [v, setV] = useState(0);
  return <Widget onReady={setV} />;
}
\end{lstlisting}

Trois \og noms\fg{} coexistent. Le nom \emph{écrit} au site d'appel (\code{Panel},
\code{Widget}) ; le couple (fichier de définition, nom exporté), qui seul désigne un composant
sans ambiguïté (\code{(b/Widget.tsx, Widget)}) ; et le nom qu'un rapport \emph{affiche}, qui
doit distinguer deux \code{Widget} (\code{Widget@b/Widget.tsx}). L'histoire du projet est
celle de la confusion entre ces trois orthographes, et de sa résolution par \adr{40}.

\subsection{\code{ComponentId} et \code{ComponentTable}}

\rustsnippet{src/ir/component_id.rs}{23}{29}{l'identité d'un composant}

Un \code{ComponentId} est un index de 4 octets, \code{Copy} et \code{Ord}. Il tient dans une
clé de store ou dans un label d'un \code{BTreeSet} sans allocation, et deux identités se
comparent en une instruction. Seule \code{ComponentTable::intern} en fabrique en production
(\code{from_index} est réservé aux tests, \src{src/ir/component_id.rs}{44}{47}). Une valeur
réservée, \code{ComponentId::SYNTHETIC} (\code{u32::MAX}, \srcline{src/ir/component_id.rs}{38}),
sert aux IR construites à la main par les tests unitaires, qui n'analysent qu'un composant et
n'en comparent jamais deux.

\rustsnippet{src/ir/component_id.rs}{61}{74}{la table d'interning}

La table est une bijection \code{origins} / \code{by_origin} entre identités et
\code{CompOrigin { file, name }}. \code{intern} est idempotent
(\src{src/ir/component_id.rs}{77}{86}) ; \code{register} accepte un identifiant existant, pour
le seul cas qui ne peut pas passer par \code{intern} : un résultat analysé seul, qui porte
\code{SYNTHETIC} dans ses labels et doit rejoindre un programme sans réécriture
(\src{src/ir/component_id.rs}{94}{101}). \code{ids()} rend les identifiants triés, et comme
l'interning suit l'ordre trié des clés du registre, tout parcours de la table est
reproductible d'une exécution à l'autre (\src{src/ir/component_id.rs}{121}{128}).

C'est le registre qui interne, parce que c'est le seul endroit qui connaît l'ensemble des
composants :

\rustsnippet{src/engine/component_registry.rs}{42}{55}{construction du registre, interning en ordre de clé trié}

La table est ensuite clonée dans le résultat de programme (champ \code{component_table} de
\code{ProgramAnalysisResult}), de sorte que règles et rendu résolvent contre
la table même qui a indexé l'analyse. Le registre lui-même est une \code{KeyedRegistry},
une table \code{(PathBuf, Symbol)} $\to$ \code{ComponentIR} partagée avec les registres de hooks
et d'utilitaires (\src{src/registry/keyed.rs}{21}{113}).

\subsection{Le nom d'affichage n'est qu'un rendu}

\rustsnippet{src/ir/component_id.rs}{138}{154}{le nom affiché, suffixé seulement en cas de collision}

Le nom d'affichage dépend du \emph{contenu} du programme analysé : \code{Widget} devient
\code{Widget@a/Widget.tsx} dès qu'un fichier sans rapport définit un second \code{Widget}.
C'est voulu --- le suffixe n'existe que pour distinguer deux homonymes dans un rapport ---, et
c'est précisément pourquoi aucune table ne doit être indexée par lui.

\begin{decision}[ADR-40 --- l'identité est un id interné, le nom d'affichage un rendu]
Avant \adr{40}, le nom d'affichage était l'orthographe \emph{sur laquelle tournait l'analyse}
(\adr{38} §5 l'avait choisie pour unifier deux orthographes) : la table des résultats, le store
partagé, le label d'un \code{Versioned}, le propriétaire d'un setter, le graphe d'appel et
\code{RuleCtx} étaient tous indexés par lui. Ajouter un fichier sans rapport renommait donc un
composant et ré-indexait toutes ces tables à la fois. \adr{40} décide : (1) l'identité est un
\code{ComponentId} interné, comme \code{FileId} l'est pour les chemins (\adr{19}) ; (2) le nom
d'affichage est minté au rendu, depuis la table, et rien ne le compare ; (3) un élément JSX
porte ce que son fichier a prouvé (\code{CompOrigin}) ; (4) la résolution vit en un seul
endroit et refuse de deviner ; (5) tout chemin est normalisé au lowering
(\code{lower_files_with}), pour que clés du registre et réponses du résolveur s'écrivent
pareil. Mesure : digest identique sur 35\,541 fichiers (1\,348 localisations), 873\,s contre
858\,s --- ni régression, ni gain revendiqué.
\end{decision}

Parmi les six tests du module (\src{src/ir/component_id.rs}{184}{274}), celui qui \og dit
pourquoi pas\fg{} s'appelle
\begin{center}\code{interning_a_namesake_changes_an_existing_display_name_but_not_its_id}\end{center}
\noindent et son nom dit tout : ajouter un homonyme
change le nom d'affichage d'un composant existant, jamais son identifiant. La couche règles
dit les deux moitiés à voix haute :

\rustsnippet[label={lst:inter-composants-rulectx-id}]{src/rules/api/query.rs}{393}{404}{l'identité pour comparer, le nom pour écrire}

\begin{piege}
Écrire \code{if ctx.component_name() == "Foo"} dans une règle est une faute, même quand un
test passe : sur un autre ensemble de fichiers, le même composant s'appelle
\code{Foo@src/a/Foo.tsx}. Toute comparaison ou toute clé passe par \code{component()} et par
la table (\code{ProgramAnalysisResult::component_named} est l'inverse public du nom
d'affichage, \src{src/engine/program_result.rs}{79}{81}, et répond \code{None} pour un nom nu
ambigu).
\end{piege}

\subsection{L'origine prouvée d'un élément JSX}

Un élément de composant porte, depuis \adr{40}, ce que \emph{son propre fichier} a prouvé du
nom qu'il écrit :

\rustsnippet{src/ir/expr.rs}{245}{258}{l'élément de composant dans l'IR}

\rustsnippet{src/ir/expr.rs}{160}{172}{l'origine : fichier de définition et nom exporté}

Les deux moitiés servent. Le fichier départage vingt composants \code{Form} ; le nom résout un
alias : \code{import { Widget as Panel }} écrit \code{<Panel/>}, mais l'origine vaut
\code{(b/Widget.tsx, Widget)}. L'origine est calculée au lowering par \code{build_jsx_origins}
(\src{src/lowering/import_resolution.rs}{237}{262}) : chaque déclaration de premier niveau du
fichier courant s'associe à elle-même, puis chaque import que le résolveur mappe vers un vrai
fichier s'associe à \code{(fichier résolu, nom importé)}. Un nom absent de la carte est
\emph{non résolu} --- import d'espace de noms, barrel de ré-export, paquet npm --- et
\code{origin} vaut \code{None}, qui se lit \og ignorance\fg{}, jamais \og absence prouvée\fg.

\subsection{\code{resolve_child} : trois réponses}

\rustsnippet{src/engine/component_registry.rs}{8}{20}{la réponse du registre sur un nom JSX}

\rustsnippet[label={lst:inter-composants-resolve-child}]{src/engine/component_registry.rs}{114}{127}{résoudre un élément, sans jamais deviner}

L'ordre est : (1) l'origine prouvée, si le registre contient ce couple ; (2) sinon le nom écrit,
s'il est unique dans tout le registre ; (3) sinon \code{Ambiguous} si plusieurs fichiers le
définissent, \code{Unknown} si aucun. Une origine qui pointe vers un fichier qui ne définit pas
le composant (un barrel, limite d'un niveau de ré-export, \issue{49}) retombe sur le nom. Le
balayage des clés est linéaire par appel.

\begin{exemple}{Alias, collision, ambiguïté}{inter-composants-amb}
Sur le répertoire de la \cref{lst:inter-composants-amb} :
\begin{sortie}
$ reactant check . --info --verbose
[verbose] 3 components analyzed
[verbose] cache hits: 3, misses: 1
  [verbose] App: 1 iteration(s), widened: []
  [verbose] Page: 1 iteration(s), widened: []
  [verbose] Widget@a/Widget.tsx: 3 iteration(s), widened: [0]
  [verbose] Widget@b/Widget.tsx: 0 iteration(s), widened: []
  Page  (1 hooks)  Page.tsx
    info   analysis-limit  several analysed files define a component called `Widget` and this reference does not resolve to one of them, so the child is treated as unknown. Import it explicitly from its file (FN possible)
    suspended  analysis-limit  4 passing check(s) withheld: the analysis was truncated in this component, so they are not guaranteed
  Widget@a/Widget.tsx  (2 hooks)  a/Widget.tsx
    warn   infinite-loop  [hook:0]  (line 5:2)  this effect keeps pushing state `n` to new values on every run. Potential infinite render loop
    …
  Widget@b/Widget.tsx  (0 hooks)  b/Widget.tsx
    error  cross-setter-in-render  var:onReady  (line 2:2)  prop `onReady` (a state setter of parent `App`) called during render of `Widget@b/Widget.tsx`, which triggers a parent re-render on every render
\end{sortie}
(sortie élaguée : la ligne du graphe de symboles, les traces et les lignes \code{verified} sont
omises.) \code{<Panel>} dans \code{App} a pour origine \code{(b/Widget.tsx, Widget)} : l'alias
est résolu, l'enfant est inliné avec un setter d'\code{App}, et l'\Error{} de
\code{b/Widget.tsx} n'existe que grâce à cet inlining. \code{<Widget>} dans \code{Page} n'a pas
d'origine (\code{Page.tsx} n'importe rien) et deux fichiers définissent le nom : réponse
\code{Ambiguous}, enfant inanalysable, \Info{} sur \code{Page}. Les noms affichés sont suffixés,
et ne servent qu'à l'affichage.
\end{exemple}

\begin{decision}[ADR-40 §4 --- refuser de deviner]
L'ancien \code{get_by_name} répondait par la première clé \code{(file, name)} \emph{dans l'ordre
de tri}. Ce n'était pas non déterministe, c'était \emph{faux} : un leurre \code{a/Widget.tsx}
remplaçait le \code{./b/Widget} importé, et l'\Error{} qui dépendait du vrai corps disparaissait
(\og ✓ no issues found\fg). Désormais \code{Ambiguous} rend l'enfant inanalysable, comme un
composant hors de l'exécution, et le parent porte une \Info{} \regle{analysis-limit}.
\textbf{Alternative refusée} : analyser chaque candidat d'un nom ambigu. C'est sound et plus
précis, et cela a été écarté sur mesure : 1\,347 références ambiguës sur les quatorze dépôts du
corpus (huit n'en ont aucune), contre 24\,500 références inconnues déjà rapportées (+5,5\,\%) ;
\code{<Button/>} seul compte 1\,453 sites ambigus quand le corpus est vu comme un seul arbre, et
payer une analyse d'enfant par candidat à chacun est la forme du blocage quadratique de
\issue{86}. Le résidu vient surtout d'alias \code{@workspace/*} non résolus (\issue{48}), dont
la correction supprime la cause.
\end{decision}

\begin{piege}
Le repli se fait sur \code{name}, le nom \emph{écrit} au site, et non sur \code{o.name}, le nom
exporté que porte l'origine. Un alias importé à travers un barrel
(\code{import { Widget as Panel } from "./index"}, où \file{index.ts} ne fait que
ré-exporter) cherche donc \code{Panel} et répond \code{Unknown}, même si un seul fichier définit
\code{Widget} (lecture de \src{src/engine/component_registry.rs}{115}{126} ; le test
\code{an_origin_pointing_at_no_component_falls_back_to_the_name} utilise le même nom des deux
côtés). Direction : un enfant inanalysable, donc havoc et \Info{} --- une perte de précision,
pas de soundness. Inversement, un fichier voit toujours ses propres déclarations dans sa carte
d'origines : un \code{<Child/>} défini dans le même fichier est résolu par l'origine, même si
d'autres fichiers définissent un \code{Child}. Vérifié au \snapshotCommit{} sur un répertoire à
trois fichiers (\code{Widget.tsx}, un barrel \code{index.ts} qui ré-exporte
\code{{ Widget }}, et \code{App.tsx} qui importe \code{{ Widget as Panel } from "./index"}) :
\code{--info --verbose} répond \Info{} \og component \code{`Panel`} was not found in the
analysis registry\fg{}, pas \og several analysed files define\fg, bien qu'un seul fichier
définisse \code{Widget}.
\end{piege}

% ===========================================================================
\section{Les racines de l'analyse}
\label{sec:inter-composants-racines}

L'inlining descendant a besoin d'un point de départ : les composants analysés sans parent.

\begin{definition}{Racine}{inter-composants-racine}
Une \terme{racine} est un composant analysé en tête de la phase 1, sans parent : son
environnement d'entrée est vide (\code{AbstractEnv::bottom()}), donc son paramètre props, non
lié, se lit $\Top$. Les enfants qu'elle rend sont inlinés avec les props qu'elle leur calcule.
\end{definition}

\rustsnippet{src/engine/root_detector.rs}{25}{36}{les trois stratégies}

Le driver choisit la stratégie d'après la ligne de commande : \code{--entry} l'emporte, puis
\code{--all-roots}, sinon l'heuristique (\src{src/driver/mod.rs}{282}{288}).

\subsection{L'heuristique : ce que personne ne rend}

\rustsnippet[label={lst:inter-composants-heuristic}]{src/engine/root_detector.rs}{57}{91}{la détection des racines, stratégie par stratégie}

Une racine heuristique est un composant qu'\emph{aucun} \code{CompApp} du programme ne
référence. La collecte, \code{collect_compapp_refs} (\src{src/engine/root_detector.rs}{126}{139}),
est syntaxique et sur-approximante : elle parcourt le rendu, les corps d'effets, de mémos, de
callbacks et de handlers, et descend explicitement dans les \code{FnLit} imbriqués (un callback
de \code{.map}, une render prop), que \code{for_each_child} ne traverse pas
(\src{src/engine/root_detector.rs}{145}{160}). Chaque référence garde le nom écrit \emph{et}
l'origine, parce que la même collecte sert aussi la relation des consommateurs de contexte
(\issue{115}), qui a besoin de l'autre projection.

La règle de marquage mérite d'être lue lentement. Une référence résolue exclut exactement un
composant. Une référence \emph{non} résolue exclut \emph{tous} les composants de ce nom, car
n'importe lequel peut être celui qui est rendu. Le commentaire du code dit ce qui arrivait
avant \issue{7} : marquer par nom seul laissait la cible d'un alias (\code{<Panel/>} pour
\code{Widget}) sans référence apparente ; elle était analysée une seconde fois comme racine,
avec des props $\Top$, et cette passe \emph{écrasait} le résultat précis que le parent avait
produit. Nous retrouverons ce phénomène d'écrasement au \cref{sec:inter-composants-cache}.

\begin{exemple}{Racines, récursion, composant inconnu}{inter-composants-roots}
\begin{lstlisting}[language=TypeScript,caption={\file{/tmp/ch16/roots.tsx}},label={lst:inter-composants-roots}]
import { useState } from "react";

function Leaf({ n }: { n: number }) { return <span>{n}</span>; }
function Tree({ depth }: { depth: number }) {
  return <Tree depth={depth - 1} />;
}
export function App() {
  const [n, setN] = useState(0);
  return (
    <div>
      <Leaf n={1} />
      <Leaf n={1} />
      <Tree depth={3} />
      <Unknown onPick={setN} />
    </div>
  );
}
export function Orphan() { return <Leaf n={2} />; }
\end{lstlisting}
\begin{sortie}
$ reactant check roots.tsx --info --verbose
[verbose] symbol graph: 4 nodes, topo order = [Tree@roots.tsx, Leaf@roots.tsx, Orphan@roots.tsx, App@roots.tsx]
[verbose] 2 components analyzed
[verbose] cache hits: 5, misses: 3
  [verbose] App: 1 iteration(s), widened: []
  [verbose] Leaf: 0 iteration(s), widened: []
  [verbose] Orphan: 0 iteration(s), widened: []
  [verbose] Tree: 0 iteration(s), widened: []
  App  (1 hooks)  roots.tsx
    info   analysis-limit  component `Unknown` was not found in the analysis registry. Pass its file on the command line to analyse it (FN possible)
    suspended  analysis-limit  4 passing check(s) withheld: the analysis was truncated in this component, so they are not guaranteed
  Tree  (0 hooks)  roots.tsx
    info   analysis-limit  recursive component reference `Tree` is not followed, so its props are treated as unknown; cross-component cycles are not fully analysed (FN possible)

✓  1 file(s) no issues found.
\end{sortie}
\code{Leaf} est référencé ; \code{Tree} se référence lui-même, il n'est donc pas racine ; les
racines sont \code{App} et \code{Orphan}. \og 2 components analyzed\fg{} compte les deux
racines et rien d'autre : le compteur \code{components_analyzed} n'est incrémenté que par les
deux boucles d'\code{analyze_program} (\srcline{src/engine/fixpoint.rs}{751},
\srcline{src/engine/fixpoint.rs}{791}), pas par l'inlining des enfants --- son commentaire
(\og including re-analyses due to fixpoint\fg, \srcline{src/engine/program_result.rs}{208}) est
inexact. \code{App} fait une itération de plus que nécessaire parce que \code{setN}, passé à
l'inconnu \code{<Unknown>}, est havoqué (\cref{sec:inter-composants-ascendant}). Les nombres de
succès et de ratés de cache s'expliqueront au \cref{sec:inter-composants-inlining}.
\end{exemple}

\subsection{\code{--all-roots} et \code{--entry}}

Avec \code{AllComponents}, tout composant est une racine, analysée avec des props $\Top$ si
aucun parent ne l'a inliné avant. Avec \code{Explicit}, la liste vient de l'utilisateur :
\code{explicit_matches} (\src{src/engine/root_detector.rs}{44}{50}) sélectionne, pour un nom nu
\code{Foo}, \emph{tous} les composants écrits \code{Foo}, et pour la forme qualifiée
\code{Foo@src/a/Foo.tsx} --- celle-là même que le rapport imprime en cas de collision --- un
seul, via \code{resolve_display_name}. Une seule fonction sert au calcul des racines et à la
détection des noms qui ne sélectionnent rien, pour que les deux ne puissent pas diverger :

\rustsnippet{src/engine/root_detector.rs}{101}{110}{les noms d'entrée qui ne sélectionnent rien}

Le driver en fait une erreur d'usage (\src{src/driver/mod.rs}{307}{321}) : une faute de frappe
dans \code{--entry} réduisait autrefois l'exécution, en silence, à une analyse intra, et
coûtait tous les résultats inter-composants avec un rapport par ailleurs propre.

\begin{sortie}
$ reactant check intro.tsx --entry Nope
[error] --entry: no component named `Nope`
[error] name a component defined in the analysed files, `Name@path` to pick one of several with the same name
\end{sortie}

\begin{piege}
Une racine explicite mal choisie ne produit pas d'erreur, seulement moins de résultats. Sur
l'\cref{lst:inter-composants-intro}, \code{--entry Badge} fait de \code{Badge} la racine,
analysée avec des props $\Top$ ; \code{Inbox}, non atteint, passe en phase 2 (intra) et y
havoque le setter qu'il passe à \code{<Badge>}. L'\Error{} disparaît (observé :
\og ✓ 1 file(s) no issues found\fg{}, \code{Inbox: 1 iteration(s)}). À l'inverse,
\code{--all-roots} la garde sur ce fichier, mais seulement parce que \code{Badge} (id 0) est
analysé comme racine \emph{avant} \code{Inbox} (id 1), qui l'inline ensuite et écrase le
résultat imprécis : l'ordre des identifiants décide (\cref{sec:inter-composants-cache}).
\end{piege}

\subsection{Un trou de la collecte, sans danger}

\code{collect_compapp_refs} ne lit pas les arguments d'un \code{HookEntry::Custom} (branche
\code{_ => {}}, \srcline{src/engine/root_detector.rs}{136}). Un élément passé à un hook opaque
--- \code{useModal(<Dialog onClose={setS}/>)}, \code{useModal} venant d'un paquet --- ne marque
donc pas \code{Dialog} comme référencé, et \code{Dialog} devient une racine, analysée avec des
props $\Top$. Le dossier de notes l'a relevé sur la sonde (racines \code{App} et
\code{Dialog}) ; le CLI montre les deux composants analysés hors inlining
(\file{/tmp/ch16/hookarg.tsx}, \og 2 components analyzed\fg, \og cache hits: 0, misses: 0\fg) et la
seule \Info{} sur le hook \code{useModal}. C'est une sur-approximation des racines, donc sûre ;
on la note parce que \code{SymbolGraph}, lui, parcourt ces arguments
(\src{src/engine/symbol_graph.rs}{271}{275}) : deux collectes \og syntaxiques\fg{} du même
périmètre ne voient pas les mêmes choses.

% ===========================================================================
\section{L'inlining descendant}
\label{sec:inter-composants-inlining}

\subsection{Le contexte partagé \code{InterCtx}}

Une racine est analysée par \code{analyze_component_impl} (\cref{chap:point-fixe-composant}),
avec un argument de plus : un contexte inter-composants.

\rustsnippet[label={lst:inter-composants-interctx}]{src/domains/context.rs}{59}{77}{le contexte d'une analyse inter-composants}

Tout l'état mutable partagé --- cache, store partagé, graphe d'appel, statistiques, table des
résultats --- est derrière des \code{RefCell}, elles-mêmes derrière des références partagées :
on passe \code{&InterCtx} partout sans conflit d'emprunts mutables. Deux champs sont propres à
un niveau : \code{component}, le composant en cours, et \code{call_stack}, la pile de ses
ancêtres. \code{child(c)} fabrique le contexte d'un enfant en copiant toutes les références et
en empilant le composant courant sur une pile neuve ; \code{is_recursive(id)} teste
\code{call_stack.contains(&id) || component == id} (\src{src/domains/context.rs}{80}{103} ; la
pile ne contient que les ancêtres, d'où le second disjoint, \cref{chap:statevalue-stabilite}).
Enfin \code{analyze_child} est un pointeur de fonction, \code{analyze_component_inter}
(\src{src/engine/fixpoint.rs}{62}{81}), fourni par le moteur : il casse le cycle de modules
entre le transfert (\file{domains}), qui rencontre l'élément, et le point fixe
(\file{engine}), qui sait analyser un composant.

\subsection{\code{eval_comp_app}, pas à pas}

L'évaluation d'un \code{Expr::CompApp} dans le domaine des valeurs délègue à
\code{eval_comp_app} (\src{src/domains/transfer/state_value.rs}{158}{163}). La fonction tient
en 125 lignes, que nous lisons en trois morceaux.

\rustsnippet[label={lst:inter-composants-eca-1}]{src/domains/transfer/state_value.rs}{476}{518}{sans contexte inter, puis résolution de l'enfant}

\begin{enumerate}
  \item \textbf{Pas d'\code{InterCtx}} : analyse d'un composant seul, phase 2, ou passe de
        rafraîchissement (\cref{sec:inter-composants-evaluations}). Tout enfant est alors
        inanalysable : on havoque les setters qui s'échappent dans ses props et on rend
        \code{Stable}.
  \item \textbf{Résolution} par \code{resolve_child}
        (\cref{lst:inter-composants-resolve-child}), puis \code{ComponentId}. \code{Unknown} et
        \code{Ambiguous} sont enregistrés séparément dans les statistiques --- ils demandent à
        l'utilisateur deux remèdes différents, élargir l'exécution ou désambiguïser --- et
        deviendront des \Info{} \regle{analysis-limit}
        (\src{src/rules/impls/analysis_limit_info.rs}{53}{77}) ; puis havoc et \code{Stable}.
\end{enumerate}

\rustsnippet[label={lst:inter-composants-eca-2}]{src/domains/transfer/state_value.rs}{520}{556}{récursion, props abstraites, cache}

\begin{enumerate}\setcounter{enumi}{2}
  \item \textbf{Récursion} : si l'enfant est déjà sur la pile, on compte une coupure, on note
        la paire \code{(courant, enfant)} et on rend \code{Stable} --- \emph{avant} le clonage
        de l'IR, qui coûte. Il n'y a \emph{pas} de havoc ici, contrairement aux branches
        \code{Unknown}/\code{Ambiguous} (\cref{sec:inter-composants-recursion}).
  \item \textbf{Props} : \code{eval_props_map} évalue chaque champ de l'objet props dans
        l'environnement courant du parent et rend une carte prop $\mapsto$ \code{EnvVal}, qui
        garde les sites du tas d'une prop fermeture ou objet. Sa version aplatie en
        \code{StateValue} sert de clé au cache et au graphe d'appel.
  \item \textbf{Cache} : un succès (égalité de treillis des props,
        \cref{sec:inter-composants-cache}) enregistre une arête du graphe d'appel et rend
        \code{Stable} \emph{sans ré-analyser l'enfant}.
\end{enumerate}

\rustsnippet[label={lst:inter-composants-eca-3}]{src/domains/transfer/state_value.rs}{558}{593}{l'environnement de l'enfant, son analyse, l'enregistrement}

\begin{enumerate}\setcounter{enumi}{5}
  \item \textbf{Environnement de l'enfant} : un environnement vide, un site frais
        \code{props_id = ExprId::fresh()} sous lequel on range, dans un tas neuf, l'objet des
        props (\code{HeapValue::Obj}), plus une copie des valeurs de tas des props qui ont
        des sites ; le paramètre de l'enfant est lié à ce site
        (\code{extend_loc}). C'est l'\code{AbstractObject} d'\adr{12} §9 : un accès
        \code{props.onSeen} dans l'enfant lit le tas et retrouve le setter du parent.
  \item \textbf{Analyse} : \code{analyze_child} lance une analyse \emph{complète} de l'enfant,
        avec son propre point fixe, sous le contexte fils.
  \item \textbf{Enregistrement} : le résultat est rangé dans \code{results} (en
        \emph{remplaçant} celui qui s'y trouvait) et dans le cache, et une arête est ajoutée
        au graphe d'appel. La valeur de l'élément est, dans tous les cas, \code{Stable}.
\end{enumerate}

La \cref{fig:inter-composants-sequence} résume ces échanges sous forme de diagramme de
séquence.

\begin{figure}[htbp]\centering
\begin{tikzpicture}[x=2.55cm,y=-0.62cm,font=\scriptsize,
    tete/.style={bloc,font=\scriptsize,minimum width=20mm,inner sep=3pt},
    vie/.style={draw=ra-gray,dashed},
    msg/.style={fleche,thin},
    ret/.style={fleche,thin,dashed},
    lab/.style={font=\scriptsize,fill=white,inner sep=1pt}]
  \foreach \x/\n/\t in {0/p/{rendu du parent},1/e/{\code{eval_comp_app}},2/r/{registre},3/c/{cache},4/k/{analyse de l'enfant},5/s/{store partagé}} {
    \node[tete] (\n) at (\x,0) {\t};
    \draw[vie] (\x,0.8) -- (\x,17.2);
  }
  \draw[msg] (0,1.6) -- node[lab,above] {\code{CompApp(name, props, origin)}} (1,1.6);
  \draw[msg] (1,2.8) -- node[lab,above] {\code{resolve_child}} (2,2.8);
  \draw[ret] (2,3.8) -- node[lab,above] {\code{Resolved} / \code{Unknown} / \code{Ambiguous}} (1,3.8);
  \node[etiquette,anchor=west,align=left] at (1.05,5.3) {inconnu, ambigu : havoc, \code{Stable}\\ sur la pile ? coupure, \code{Stable}\\ sinon \code{eval_props_map}};
  \draw[msg] (1,7.0) -- node[lab,above] {\code{lookup(child, props)}} (3,7.0);
  \draw[ret] (3,8.1) -- node[lab,above] {succès : arête, \code{Stable}} (1,8.1);
  \draw[msg] (1,9.4) -- node[lab,above] {\code{analyze_child(ir, env, tas, ctx fils)}} (4,9.4);
  \draw[msg] (4,10.8) -- node[lab,above] {\code{update((parent, l), v)}} (5,10.8);
  \node[etiquette,anchor=west,align=left] at (4.05,11.8) {appel d'un setter\\ du parent};
  \draw[ret] (4,13.0) -- node[lab,above] {\code{AnalysisResult}} (1,13.0);
  \draw[msg] (1,14.0) -- node[lab,above] {\code{insert} ; \code{results.insert}} (3,14.0);
  \draw[ret] (1,15.0) -- node[lab,above] {\code{Stable}} (0,15.0);
  \draw[msg] (0,16.4) -- node[lab,above] {\code{slice(parent)} au test de convergence} (5,16.4);
\end{tikzpicture}
\caption{Diagramme de séquence de l'évaluation d'un élément \code{<Child/>} sous un
\code{InterCtx}. Les retours sont en pointillés. L'écriture de l'enfant dans le store partagé
n'est lue par le parent qu'au test de convergence de \emph{sa} boucle, qui relance un tour si
elle fait monter son état.}
\label{fig:inter-composants-sequence}
\end{figure}

\subsection{Les props dans le tas}
\label{sec:inter-composants-props-tas}

Pourquoi passer les props par le tas plutôt que par l'environnement ? Parce qu'un enfant doit
pouvoir \emph{exécuter} une fermeture reçue en prop. \code{eval_props_map}
(\src{src/domains/transfer/state_value.rs}{653}{689}) alloue dans le tas du parent un
\code{FnLit} passé en ligne (\code{onTick={() => setB(b + 1)}}), avec ses captures, et garde
son site ; une prop qui aliase un littéral déjà alloué garde aussi ses sites
(\code{resolve_locs}). L'enfant reçoit une copie de ces entrées : quand son effet appelle
\code{onTick()}, le corps s'exécute avec ses captures, et l'appel de \code{setB} a pour
propriétaire le parent (\cref{chap:transfert-stores}, exemple \og une boucle croisée\fg).

Deux limites de cette frontière sont traitées ailleurs, et il suffit ici de les situer. Seuls
les sites \emph{directs} des props sont copiés : une fermeture rangée comme \emph{membre} d'un
objet passé en prop arrive dans l'enfant sans son corps (défaut D11 du dossier 05,
\cref{chap:transfert-stores}). Et les sites copiés gardent les \code{ExprId} du fichier parent,
qui peuvent entrer en collision avec ceux que l'enfant alloue dans son propre fichier
(\cref{chap:ir,chap:splice-inlining}).

\subsection{Combien de fois un élément est-il évalué ?}
\label{sec:inter-composants-evaluations}

Les compteurs de cache de l'\cref{ex:inter-composants-roots} (5 succès, 3 ratés) paraissent
énigmatiques tant qu'on ne sait pas combien de fois un \code{CompApp} est évalué. Trois faits,
vérifiés dans le code, suffisent.

\begin{enumerate}
  \item La passe de rendu d'un composant tourne \code{iterations + 1} fois sous
        \code{InterCtx} : la boucle externe sort sur \code{new_state.leq(&state)}
        \emph{avant} d'incrémenter \code{iteration} (\src{src/engine/fixpoint.rs}{355}{499}).
  \item Elle est suivie d'une passe de rafraîchissement \emph{sans} \code{InterCtx}
        (\src{src/engine/fixpoint.rs}{532}{560}), qui ne ré-inline rien : elle passe par la
        branche \og pas d'\code{InterCtx}\fg{} d'\code{eval_comp_app}.
  \item Dans une passe, un élément imbriqué dans un élément hôte est évalué \emph{une} fois
        (bras \code{CompApp} d'\code{exec_callbacks_depth},
        \src{src/domains/interp/interpreter.rs}{529}{533}) ; un élément qui est
        \emph{directement} l'expression du \code{return} l'est \emph{deux} fois : par ce même
        bras, puis par l'appel d'\code{exec_expr_effects} sur le terminateur \code{Return}
        (\src{src/domains/interp/interpreter.rs}{123}{125},
        \src{src/engine/cfg_analyzer.rs}{84}{86}). La seconde évaluation est un succès de
        cache.
\end{enumerate}

\begin{exemple}{Le \code{return} direct compte double}{inter-composants-double}
Dans \file{/tmp/ch16/ret.tsx}, \code{function A() { return <Leaf n={1} />; }} et
\code{function B() { return <div><Leaf n={2} /></div>; }} : \code{--verbose} rapporte
\og cache hits: 1, misses: 2\fg. Un raté par valeur de props, et un seul succès : la seconde
évaluation du \code{return} direct d'\code{A}. Sur \file{roots.tsx}, \code{App}
(\code{iterations = 1}, donc deux passes, éléments imbriqués dans \code{<div>}) évalue ses
quatre éléments deux fois ; \code{<Unknown>} ne touche pas le cache ; \code{Orphan}
(\code{return <Leaf n={2} />} direct) évalue son élément deux fois. On obtient trois ratés
(\code{Leaf{n:1}}, \code{Tree}, \code{Leaf{n:2}}) et cinq succès (le second
\code{<Leaf n={1}/>} de la première passe, les trois éléments résolus de la seconde passe
d'\code{App}, la seconde évaluation d'\code{Orphan}). Même raisonnement sur \file{drill.tsx} :
\og cache hits: 3, misses: 7\fg{} (\cref{exo:inter-composants-drill-cache}).
\end{exemple}

Le coût d'un élément est donc au plus une analyse de l'enfant par valeur de props distincte,
bornée par le cache, mais cette analyse est elle-même récursive (les enfants de l'enfant) : le
coût total dépend de la profondeur de l'arbre et du nombre de passes de chaque ancêtre. Il n'y
a pas de borne plus fine dans le code ; la pile interdit seulement les cycles.

\subsection{La récursion : une coupure, pas un havoc}
\label{sec:inter-composants-recursion}

Un composant récursif est légitime en React --- un arbre, un menu imbriqué --- et chaque
instance a son propre état. L'analyse, elle, ne suit pas la récursion : \adr{12} §11 prévoyait
une coupure \og façon MOPSA\fg{} avec un résultat $\Top$ et une hypothèse
\code{A_ignore_recursion} enregistrée.

\begin{exemple}{Un setter remis à l'instance intérieure}{inter-composants-rec}
\begin{lstlisting}[language=TypeScript,caption={\file{/tmp/ch16/rec.tsx}},label={lst:inter-composants-rec}]
import { useState } from "react";
function Tree({ depth, onHit }: { depth: number; onHit?: (v: number) => void }) {
  const [x, setX] = useState(0);
  if (depth < 3) onHit?.(1);
  return <div>{x}{depth > 0 ? <Tree depth={depth - 1} onHit={setX} /> : null}</div>;
}
export function App() {
  return <Tree depth={3} />;
}
\end{lstlisting}
À l'exécution, l'instance intérieure appelle \code{onHit}, c'est-à-dire le \code{setX} de
l'instance extérieure, pendant son rendu : un vrai \regle{cross-setter-in-render}.
\begin{sortie}
$ reactant check rec.tsx --info --show-clean
  Tree  (1 hooks)  rec.tsx
    info   analysis-limit  recursive component reference `Tree` is not followed, so its props are treated as unknown; cross-component cycles are not fully analysed (FN possible)
    suspended  analysis-limit  4 passing check(s) withheld: the analysis was truncated in this component, so they are not guaranteed

✓  1 file(s) no issues found.
\end{sortie}
\end{exemple}

L'analyse inline \code{Tree} une fois depuis \code{App} ; l'élément \code{<Tree …/>} intérieur
tombe sur la coupure, qui rend \code{Stable} \emph{sans} havoc : \code{x} n'est pas mis à $\Top$,
et l'écriture de l'instance intérieure n'est vue par personne. Le faux négatif est
\emph{déclaré} --- \Info{} \regle{analysis-limit} et suspension des assurances
\og verified\fg{} du composant tronqué --- et c'est cette suspension, non une
sur-approximation de l'état, qui protège les garanties publiées.

\begin{piege}
Le texte de l'\Info{} (\og its props are treated as unknown\fg) décrit l'intention d'\adr{12},
pas le code : l'instance récursive n'est pas analysée, ses props ne sont ni évaluées ni
havoquées. De même, les commentaires de \code{AnalysisStats::recursive_component_refs}
(\og cut to ⊤\fg, \srcline{src/engine/program_result.rs}{210}) et de
\code{ProgramAnalysisResult::recursive_components} (\og received ⊤ result\fg,
\srcline{src/engine/program_result.rs}{33}) sont inexacts : la coupure rend \code{Stable} et ne
met rien à $\Top$ (\src{src/domains/transfer/state_value.rs}{520}{528}). Havoquer ici, avec le
même mécanisme que la branche \code{Unknown}, serait la correction centrale ; le choix n'est
pas tranché au snapshot.
\end{piege}

% ===========================================================================
\section{Le flux ascendant : setters, store partagé, havoc}
\label{sec:inter-composants-ascendant}

\subsection{Un setter qui porte son propriétaire}

\adr{12} §7 introduisait une valeur \code{ComponentSetter { component, label }}, transposition
du \code{Set_clos { label, path }} de React-tRace~\cite{LeeAhnYi2025} : un setter qui sait à
quel composant il appartient. Elle est devenue la composante setter du produit
\code{StateValue} (\cref{def:statevalue}), dont \code{as_setter()} rend le couple
\code{(ComponentId, HookLabel)}. Quand un parent passe \code{setSeen}, la prop \code{onSeen} de
l'enfant vaut ce setter, propriétaire compris. L'interpréteur n'a plus qu'à regarder, à chaque
appel, si le callee s'évalue en setter :

\rustsnippet[label={lst:inter-composants-cross-setter}]{src/domains/interp/interpreter.rs}{368}{391}{l'appel d'un setter porteur de propriétaire écrit le store partagé}

L'argument (le premier, ou $\Top$ s'il manque ou ne se lit pas) est joint dans l'entrée
\code{(propriétaire, label)} du store partagé. Le bras est gardé par \code{ctx.inter.is_some()} :
sans contexte inter, il ne fait rien --- mais sans contexte inter, aucun setter étranger ne
peut non plus arriver dans l'environnement.

\subsection{Le store partagé, et pourquoi il n'y a pas de point fixe de programme}

\rustsnippet{src/domains/stores/shared_state_store.rs}{29}{34}{une mise à jour monotone}

Le store partagé (\cref{def:store}) associe à chaque couple \code{(ComponentId, HookLabel)} une
valeur, $\Bot$ pour une clé inconnue, et ne fait que croître. La boucle externe du parent
importe sa tranche à chaque test de convergence :

\rustsnippet[label={lst:inter-composants-import}]{src/engine/fixpoint.rs}{486}{497}{le parent importe les écritures de ses enfants}

\begin{decision}[ADR-12 §1, §5, §8 --- inlining descendant, flux bidirectionnel, pas de couche programme]
Le parent est analysé d'abord, et l'enfant est analysé au point d'application avec les props
abstraites ; pas de résumés paramétriques ascendants, qui demanderaient une couche d'abstraction
supplémentaire \og without justified gain at the current stage\fg. Le flux est bidirectionnel :
les props descendent, les setters remontent. Et il n'y a \og no separate program-level fixpoint
layer\fg{} : les écritures d'un enfant sont des mutations du store partagé,
\og indistinguishable from local setters from the fixpoint's perspective\fg. La sensibilité au
contexte est obtenue par le cache (une analyse par valeur de props). Écarts avec le texte :
les clés sont des \code{ComponentId}, pas des \code{Symbol} (\adr{40}) ; la récursion rend
\code{Stable}, pas $\Top$.
\end{decision}

\begin{proposition}{L'import d'un store monotone préserve la convergence}{inter-composants-monotone}
Soit $S_i$ l'état du parent au tour $i$ et $X_i$ sa tranche du store partagé lue à ce tour. Si
$X_i \lleq X_{i+1}$ pour tout $i$ (le store ne fait que croître) et si le nouvel état est
$S'_i = F(S_i) \join X_i$ avec $F$ monotone, la boucle externe du parent, munie de son widening
(\cref{chap:point-fixe-composant}), s'arrête sur un état $S$ tel que $F(S) \join X \lleq S$,
où $X$ est la tranche lue au dernier tour.
\end{proposition}

\noindent\emph{Esquisse.} La suite des tranches est croissante dans un domaine où le widening
force la stabilisation des slots ; le test \code{new_state.leq(&state)} est le même que pour un
setter local, et une écriture d'enfant n'est, pour la boucle, qu'un terme de plus dans le join.
La proposition ne dit rien de la tranche écrite \emph{après} le dernier tour du parent (par un
enfant analysé plus tard, sous une autre racine) : c'est l'une des raisons pour lesquelles
la table des résultats n'est pas un point fixe de programme
(\cref{sec:inter-composants-cache}). \hfill$\square$

Sur l'\cref{lst:inter-composants-intro}, la chaîne est complète : \code{Inbox} rend
\code{<Badge>}, qui est inliné ; le rendu de \code{Badge} appelle \code{onSeen(1)}, setter
d'\code{Inbox} : \code{shared[(Inbox, 0)]} $\join$= $[1,1]$. Au test de convergence,
\code{Inbox} importe sa tranche, son état monte, il fait un second tour (\og Inbox: 1
iteration(s)\fg{} en \code{--verbose}), où \code{<Badge>} est un succès de cache (la prop
\code{onSeen} n'a pas changé : \og cache hits: 1, misses: 1\fg). La règle
\regle{cross-setter-in-render}, qui tourne sur \code{Badge}, voit dans son résultat un appel
inconditionnel, pendant le rendu, d'une valeur setter dont le propriétaire est un autre
composant.

\subsection{Le havoc des setters qui s'échappent}

\begin{definition}{Havoc}{inter-composants-havoc}
\terme{Havoquer} un slot, c'est y joindre $\Top$ parce que son setter s'échappe vers un
code que l'analyse ne peut pas suivre : un enfant inconnu, ambigu, ou toute analyse sans
contexte inter. Un tel code peut appeler le setter n'importe quand, avec n'importe quel
argument ; laisser le slot intact sous-approximerait l'état et fabriquerait des conclusions
\og l'état est stable\fg.
\end{definition}

\code{havoc_setter_props} (\src{src/domains/transfer/state_value.rs}{265}{297}, détaillé au
\cref{chap:transfert-stores}) collecte, pour chaque champ de l'objet props, le setter nu
(\code{as_setter()} sur la valeur) et les setters cachés derrière des fermetures, des objets du
tas et des spreads (\code{collect_escaping_setters}). Puis :

\rustsnippet{src/domains/transfer/state_value.rs}{287}{296}{un setter propre au composant ou d'un ancêtre}

Un setter du composant courant est havoqué dans son propre store ; un setter d'ancêtre, arrivé
par une prop et transmis plus bas, l'est dans le store partagé, que l'ancêtre importera.
Sur \file{roots.tsx}, c'est ce havoc de \code{setN}, remis à \code{<Unknown>}, qui fait faire
à \code{App} un second tour.

\begin{exemple}{Un setter dans un littéral inline n'est pas havoqué}{inter-composants-hv}
\begin{lstlisting}[language=TypeScript,caption={\file{/tmp/ch16/hv.tsx}},label={lst:inter-composants-hv}]
import { useState } from "react";
export function A() { const [n, setN] = useState(0); return <div>{n}<Unknown onPick={setN} /></div>; }
export function B() { const [n, setN] = useState(0); return <div>{n}<Unknown api={{ setN }} /></div>; }
export function C() { const [n, setN] = useState(0); const api = { setN }; return <div>{n}<Unknown api={api} /></div>; }
export function D() { const [n, setN] = useState(0); return <div>{n}<Unknown list={[setN]} /></div>; }
\end{lstlisting}
\begin{sortie}
$ reactant check hv.tsx --verbose
  [verbose] A: 1 iteration(s), widened: []
  [verbose] B: 0 iteration(s), widened: []
  [verbose] C: 1 iteration(s), widened: []
  [verbose] D: 0 iteration(s), widened: []
\end{sortie}
(lignes \code{[verbose]} des composants seulement.) Le havoc, qui fait bouger \code{n} au premier
tour et donc provoque un second tour, a lieu pour \code{A} (setter nu) et pour \code{C} (objet
du tas), \emph{pas} pour \code{B} ni \code{D}. \code{collect_escaping_setters}
(\src{src/domains/transfer/state_value.rs}{315}{367}) descend dans les champs d'un
\code{ObjectLit} ou d'un \code{ArrayLit} syntaxique, mais n'y retient un champ \code{Var} que
s'il désigne un site du tas ; un setter nu en position de champ n'est poussé nulle part, et
\code{havoc_setter_props} ne teste \code{as_setter()} que sur la valeur du champ de props
entier. L'état reste à sa valeur initiale : c'est un trou de soundness du domaine des valeurs,
dont le cas le plus courant est \code{<Ctx.Provider value={{ q, setQ }}>} --- un provider
est un composant inconnu du registre (\cref{sec:inter-composants-contextes}). Le dossier de
notes n'a construit aucun diagnostic faux à partir de ce trou ; il n'avait pas d'issue au
snapshot.
\end{exemple}

% ===========================================================================
\section{Le cache de composants}
\label{sec:inter-composants-cache}

\subsection{Une analyse par valeur de props}

\rustsnippet{src/engine/component_cache.rs}{19}{27}{le cache, indexé par identité et props abstraites}

La clé est l'identité de l'enfant et ses props abstraites aplaties en \code{StateValue} (les
sites du tas sont perdus à l'aplatissement). Chaque composant a au plus
\code{DEFAULT_MAX_PER_COMPONENT = 5} entrées (\srcline{src/engine/component_cache.rs}{10}).

\rustsnippet{src/engine/component_cache.rs}{98}{119}{égalité stricte : mêmes clés, $\lleq$ dans les deux sens}

Le succès exige l'égalité de treillis, $a \lleq b \wedge b \lleq a$, champ par champ, après une
comparaison des cardinaux. \adr{12} §2 l'explique : ne jamais réutiliser un résultat calculé
pour une entrée moins précise (par exemple $\Top$) quand l'entrée courante est plus précise.

\rustsnippet{src/engine/component_cache.rs}{65}{92}{insertion, et éviction par join}

Au-delà de cinq entrées, toutes les props sont jointes point à point, une clé absente d'une
entrée comptant pour $\Top$ (\code{join_all_props}, \src{src/engine/component_cache.rs}{121}{136}),
en une seule entrée dégradée.

\begin{exemple}{Égalité stricte et éviction}{inter-composants-cache}
\file{/tmp/ch16/cache1.tsx} (\code{<Leaf n={1}/><Leaf n={1}/>} sous un parent sans état) :
\og cache hits: 1, misses: 1\fg{} ; le second site est un succès. \file{/tmp/ch16/cache.tsx}
(\code{n} vaut 1 à 6, puis encore 1) : \og cache hits: 0, misses: 7\fg. Au sixième insert, les
cinq entrées et la nouvelle sont jointes en \code{n} $\in [1,6]$ ; le septième site
(\code{n = 1}) rate, car $[1,1] \neq [1,6]$, et il est ré-inséré. Le test unitaire
\code{eviction_on_overflow} (\src{src/engine/component_cache.rs}{281}{299}) vérifie qu'il ne
reste qu'une entrée après débordement ; \code{top_props_match_anything_after_eviction}, malgré
son nom, vérifie qu'un \code{Stable} \emph{rate} contre une entrée $\Top$
(\src{src/engine/component_cache.rs}{301}{326}).
\end{exemple}

\subsection{Un \og déjà vu\fg{}, pas une table de résultats}

Voici la subtilité la plus importante du chapitre. \code{lookup} renvoie un
\code{Arc<AnalysisResult>}, mais son seul appelant, \code{eval_comp_app}, ne regarde que
\code{is_some()} (\cref{lst:inter-composants-eca-2}). Trois conséquences :
\begin{itemize}
  \item le résultat stocké dans une entrée n'est \emph{jamais relu} ;
  \item le résultat \og de programme\fg{} d'un enfant est \code{results[child]}, remplacé à
        chaque raté par l'analyse du dernier site, jamais joint ;
  \item l'entrée dégradée associe des props jointes au résultat de la \emph{dernière}
        insertion : la doc (\og sound over-approximation\fg) ne vaut que pour la décision
        \og sauter une analyse\fg.
\end{itemize}

Sauter une analyse est sound pour le \emph{store partagé} : l'analyse sautée a déjà fait ses
écritures, de façon monotone, lors du raté qui a rempli l'entrée. Elle ne l'est pas pour la
table des résultats, que lisent toutes les règles intra-composant de l'enfant.

\begin{exemple}{Un résultat d'enfant écrasé : un faux négatif qui dépend de l'ordre}{inter-composants-overwrite}
\begin{lstlisting}[language=TypeScript,caption={\file{/tmp/ch16/overwrite.tsx}},label={lst:inter-composants-overwrite}]
import { useState } from "react";

function Child({ onUpdate }: { onUpdate: (n: number) => void }) {
  onUpdate(1);
  return <div />;
}
export function A() {
  const [n, setN] = useState(0);
  return <Child onUpdate={setN} />;
}
export function B() {
  return <Child onUpdate={() => {}} />;
}
\end{lstlisting}
\begin{sortie}
$ reactant check overwrite.tsx
   2 clean component(s) hidden, rerun with --show-clean

✓  1 file(s) no issues found.
\end{sortie}
Renommons \code{A} en \code{Z}, sans rien changer d'autre (\file{/tmp/ch16/overwrite2.tsx}) :
\begin{sortie}
$ reactant check overwrite2.tsx
  Child  (0 hooks)  overwrite2.tsx
    error  cross-setter-in-render  var:onUpdate  (line 4:2)  prop `onUpdate` (a state setter of parent `Z`) called during render of `Child`, which triggers a parent re-render on every render
\end{sortie}
(dernières lignes omises.) Les racines sont analysées dans l'ordre des identifiants, donc
alphabétique dans un même fichier. Dans le premier cas, \code{A} inline \code{Child} avec un
setter, puis \code{B} l'inline avec une fonction vide --- props différentes, raté --- et
\emph{remplace} \code{results[Child]}. La règle ne voit plus que le dernier site. Dans le
second, \code{Z} passe en dernier et son site gagne.
\end{exemple}

L'erreur réelle d'\code{A} est perdue, sans \Info{} ni suspension : c'est un faux négatif non
déclaré, qui viole l'invariant de soundness du projet. Il n'avait pas d'issue au snapshot (le
dossier de notes l'a signalé). La même famille couvre \code{--all-roots} : chaque enfant y est
ré-analysé comme racine, avec des props $\Top$, et si son identifiant est plus grand que celui
de son parent, ce dernier résultat l'emporte (la boucle de phase 1 ne saute pas une racine déjà
présente dans \code{results}). La correction centrale n'est pas triviale : joindre les résultats
de tous les sites, ou indexer les résultats par (identité, props) et faire lire aux règles
l'ensemble des contextes d'un composant (\cref{exo:inter-composants-overwrite}).

\begin{piege}
Deux résultats d'un même composant ne se joignent pas comme deux états : un
\code{AnalysisResult} porte des CFG, des tables de hooks, des relations, des environnements par
bloc. C'est pourquoi \code{exit_env} lui-même joint avec \code{reduce} et non avec
\code{fold(⊥, join)} : dans ce treillis d'environnements, $\Bot$ n'est pas neutre pour le join,
une clé absente de $\Bot$ devenant $\Top$ (\src{src/engine/analysis_result.rs}{275}{288}).
\end{piege}

% ===========================================================================
\section{Deux phases}
\label{sec:inter-composants-phases}

\rustsnippet[label={lst:inter-composants-phase1}]{src/engine/fixpoint.rs}{720}{755}{phase 1 : les racines, sous un \code{InterCtx}}

La phase 1 analyse chaque racine, dans l'ordre croissant des identifiants, avec un contexte
inter neuf (pile vide) ; ses enfants sont inlinés à la volée, et leurs résultats s'accumulent
dans \code{results}. La phase 2 (\src{src/engine/fixpoint.rs}{764}{794}, commentée au
\cref{chap:point-fixe-composant}) balaie tous les composants qui ne sont ni racines ni présents
dans \code{results}, et les analyse \emph{intra} : \code{inter = None}, props $\Top$, mais avec
leur vrai \code{ComponentId}, parce qu'un label \code{Versioned} minté ici sera comparé à des
résultats de phase 1. Entre les deux, une photographie :

\rustsnippet{src/engine/fixpoint.rs}{757}{762}{ce que la phase 1 a atteint, photographié avant la phase 2}

\begin{figure}[htbp]\centering
\begin{tikzpicture}[node distance=8mm and 14mm]
  \node[bloc,very thick] (app) {\code{App}\\{\scriptsize racine}};
  \node[bloc,very thick,right=30mm of app] (page) {\code{Page}\\{\scriptsize racine}};
  \node[bloc,below=of app] (wb) {\code{Widget@b/Widget.tsx}\\{\scriptsize inliné, props d'\code{App}}};
  \node[bloc,draw=ra-gray,fill=white,right=22mm of page] (wa) {\code{Widget@a/Widget.tsx}\\{\scriptsize phase 2 : intra, props $\Top$}};
  \node[cfgbloc,below=of page] (amb) {\code{<Widget>} ambigu\\havoc de \code{setV}};
  \draw[flechemust] (app) -- node[etiquette,left] {\code{<Panel>}} (wb);
  \draw[flechemay] (page) -- (amb);
  \begin{scope}[on background layer]
    \node[draw=ra-blue,dashed,rounded corners,fill=ra-bg!40,inner sep=3mm,fit=(app)(page)(wb)(amb),label={[etiquette]above:\code{phase1_reached}}] {};
  \end{scope}
  \node[etiquette,below=2mm of wa,align=center] {aucune arête du graphe d'appel ;\\ \code{callers_of} vide = ascendance inconnue};
\end{tikzpicture}
\caption{Les deux phases sur l'exemple de la \cref{lst:inter-composants-amb}. En phase 1, les
racines \code{App} et \code{Page} sont analysées ; \code{App} inline \code{Widget@b} par son
origine ; la référence nue de \code{Page} est ambiguë, donc havoquée. Personne n'atteint
\code{Widget@a}, que la phase 2 analyse seul. L'ensemble \code{phase1_reached} (cadre
pointillé) est photographié entre les deux.}
\label{fig:inter-composants-phases}
\end{figure}

\subsection{La discipline de lecture de \code{phase1_reached}}

Un composant de phase 2 n'a pas de \code{InterCtx}, donc n'enregistre aucune arête du graphe
d'appel : vu par \code{callers_of}, il est indiscernable d'une vraie racine. C'est tout le sens
de la photographie (\issue{110}), dont la documentation énonce la règle :

\rustsnippet{src/engine/program_result.rs}{52}{68}{ce que la phase 1 a atteint, et comment le lire}

\begin{invariant}{Une ascendance vide n'est pas une ascendance prouvée}{inter-composants-ascendance}
Pour un composant hors de \code{phase1_reached}, un \code{callers_of} vide signifie
\og ascendance inconnue\fg, jamais \og racine prouvée\fg. Toute relation qui remonte les appelants
pour conclure à une \emph{absence} --- pas de provider au-dessus, pas d'appelant qui détienne le
setter --- consulte d'abord cet ensemble.
\end{invariant}

\code{complete_ancestry} (\src{src/engine/program_result.rs}{134}{151}) rend la fermeture
transitive des appelants, ou \code{None} dès qu'un composant du chemin n'a pas été analysé en
phase 1. Son consommateur est la relation des consommateurs de contexte
(\srcline{src/rules/helpers/context_flow.rs}{100}), qui écarte aussi une ascendance passant par
une récursion coupée (\code{recursive_components},
\src{src/rules/helpers/context_flow.rs}{104}{108}). La limite ouverte \issue{20} en est la
conséquence visible : \regle{cross-component-infinite-loop} reste muette quand le parent n'est
analysé qu'en intra.

Le \cref{chap:point-fixe-composant} a listé ce que \code{inter = None} coûte en phase 2 : hooks
personnalisés non greffés, \code{SummaryRegistry} non consulté, tout enfant havoqué, troncature
du budget d'utilitaires silencieuse. Ajoutons-en une, symétrique de la dernière :
\code{callback_depth_capped} (\src{src/domains/interp/interpreter.rs}{484}{491}) et
\code{inline_budget_exhausted} (\src{src/engine/fixpoint.rs}{211}{218}) ne sont renseignés que
sous \code{InterCtx}, si bien qu'un composant de phase 2 qui atteint ces plafonds ne porte pas
l'\Info{} correspondante. Les résumés de dépendance de rendu, eux, restent précis en phase 2,
parce qu'ils sont locaux (\cref{sec:inter-composants-render-deps}).

% ===========================================================================
\section{Ce que lit une règle}
\label{sec:inter-composants-resultats}

\subsection{\code{AnalysisResult} : le résultat d'un composant}

Le résultat d'un composant (\src{src/engine/analysis_result.rs}{171}{272}) est ce que lisent
toutes les règles. Ses champs se rangent en cinq groupes (\cref{tab:inter-composants-result}) ;
les chapitres précédents ont décrit comment chacun est calculé.

\begin{table}[htbp]\centering\small
\begin{tabular}{p{27mm}p{50mm}p{62mm}}
\toprule
groupe & champs & rôle \\
\midrule
identité & \code{component}, \code{file}, \code{param}, \code{dom_props}, \code{module_consts} &
  le composant (son \code{ComponentId}), son fichier (clé des témoins, \adr{19}), le binding
  des props, les props typées DOM, les constantes de module (seule preuve qu'un
  \code{<X.Provider>} est un provider) \\
stores convergés & \code{state_store}, \code{memo_store}, \code{block_states},
  \code{effect_block_states}, \code{handler_block_states}, \code{heap},
  \code{effect_setter_writes}, \code{custom_arg_returns} &
  valeurs des slots et des mémos, environnement en sortie de chaque bloc (rendu, effets,
  handlers), tas final, écritures pures des effets (\cref{chap:point-fixe-composant}) \\
tables syntaxiques & \code{render_cfg}, \code{hooks}, \code{hook_provenance},
  \code{hook_calls}, \code{effect_info}, \code{handler_info} &
  le CFG après expansion, les entrées de hooks et leur provenance, les sites d'appel, les deps,
  les handlers \\
relations & \code{slot_writers}, \code{slot_seeds}, \code{registrations},
  \code{effect_triggers} & calculées à convergence (\cref{chap:relations}) \\
traçabilité & \code{widen_trace}, \code{inline_origins}, \code{iterations} &
  widenings forcés, symboles inlinés, nombre d'itérations externes \\
\bottomrule
\end{tabular}
\caption{Les champs d'\code{AnalysisResult}. Tout champ documenté \og Empty for hand-built
IR\fg{} se lit \og non prouvé\fg, jamais \og absence prouvée\fg.}
\label{tab:inter-composants-result}
\end{table}

Un seul point est propre à ce chapitre : le résultat d'un enfant est celui de sa \emph{dernière}
analyse effective, pour les props de son dernier site analysé
(\cref{ex:inter-composants-overwrite}).

\subsection{\code{ProgramAnalysisResult} et ses satellites}

Le résultat de programme (\src{src/engine/program_result.rs}{28}{69}) rassemble : la table
\code{components} (\code{ComponentId} $\to$ \code{AnalysisResult}) ; le store partagé final ; le
graphe d'appel ; les composants dont la récursion a été coupée ; les statistiques ; la
\code{ComponentTable} ; la \code{FileTable} (\adr{19}) ; le registre des utilitaires, exposé aux
producteurs de témoins ; la \code{ModuleTable} (\adr{26}) ; et \code{phase1_reached}.
\code{file_table} et \code{module_table} sont laissés vides par \code{analyze_program} et remplis
par \code{analyze_lowered} (\src{src/resolver/mod.rs}{448}{451}). \code{Default} est \og le
programme vide\fg, où chaque table répond \og rien de connu\fg ; \code{single}
(\src{src/engine/program_result.rs}{91}{114}) fabrique un programme d'un seul composant déjà
analysé en \emph{enregistrant} son identifiant (souvent \code{SYNTHETIC}) au lieu d'en minter
un que ses labels internes ne connaîtraient pas.

Le graphe d'appel (\src{src/engine/program_result.rs}{166}{201}) associe à chaque appelant une
liste de \code{CallSite { callee, props, location }}. Il est bruité, et il faut le savoir pour
le lire : \code{location} vaut toujours \code{None} (\code{record_call_site(..., None)}), une
ligne est poussée à \emph{chaque} évaluation --- donc plusieurs par passe de rendu, et deux
pour un \code{return} direct (\cref{sec:inter-composants-evaluations}) ---, et
\code{callers_of} balaie toutes les arêtes et rend chaque appelant une fois, alors que
\code{callees_of} rend les lignes telles quelles, doublons compris.
L'information par site, avec sa position, est aujourd'hui portée par les \code{ElementSite} de
la dépendance de rendu (\cref{sec:inter-composants-collecte}).

Les statistiques \code{AnalysisStats} (\src{src/engine/program_result.rs}{203}{225}) comptent
succès et ratés du cache, coupures de récursion et composants analysés, et rangent en ensembles
les paires \code{(appelant, appelé)} coupées par récursion, les références inconnues et
ambiguës (le callee y reste un nom, faute de résolution), les composants dont la profondeur de
callbacks ou le budget d'utilitaires a été atteint. La règle \regle{analysis-limit}
(\src{src/rules/impls/analysis_limit_info.rs}{39}{97}) les transforme en \Info{} et suspend les
assurances du composant concerné.

\subsection{\code{RuleCtx} et \code{ProgramCache}}

Une règle ne reçoit ni l'un ni l'autre directement, mais un \code{RuleCtx}
(\src{src/rules/api/query.rs}{318}{323}) : l'identité du composant, son \code{AnalysisResult}
résolu une fois, sa configuration, et un \code{ProgramCache} partagé par tous les composants du
programme. \code{ctx.comp()} rend le résultat du composant, \code{ctx.program()} celui du
programme, \code{ctx.component()} l'identité, \code{ctx.component_name()} le nom à écrire
(\cref{lst:inter-composants-rulectx-id}). Le \code{ProgramCache}
(\src{src/rules/api/cache.rs}{26}{31}) compose les relations de programme du moteur
(\code{ProgramRelations}, le graphe de churn, \adr{42} §5) avec trois structures que la couche
règles construit encore elle-même, chacune dans un \code{OnceLock} calculé à la première
demande : la relation des consommateurs de contexte, l'index des montages et l'index de rendu
(\code{RenderIndex}). Demander les montages construit l'index de rendu, puisque
\code{MountIndex} lit ses conditions de montage sur \code{ElementSite::guard}
(\src{src/rules/api/cache.rs}{62}{68}, \issue{149}).

% ===========================================================================
\section{Le graphe de symboles : prévu, réalisé}
\label{sec:inter-composants-symboles}

\adr{13} (analyse inter-fichiers) prévoyait un pipeline \emph{eager} en cinq étapes, dont la
dernière était \og Lower + analyze in topo order\fg{} : un graphe de symboles (composants, hooks,
utilitaires --- pas de fichiers, parce que les imports circulaires entre fichiers sont courants
alors que les dépendances circulaires entre fonctions React sont quasi inexistantes), trié
topologiquement pour analyser les feuilles d'abord. Le graphe existe
(\file{src/engine/symbol_graph.rs}) : \code{build} (\src{src/engine/symbol_graph.rs}{60}{118})
crée un nœud par composant et par hook et une arête $A \to B$ quand $A$ appelle syntaxiquement
$B$ --- précise quand le lowering a résolu le callee (origine d'un élément,
\code{resolved_file} d'un hook), sinon biaisée vers le même fichier puis vers le premier nœud
homonyme inséré. \code{topo_sort} (\src{src/engine/symbol_graph.rs}{195}{255}) est un Kahn
inversé : les appelés d'abord, les nœuds pris dans des cycles ajoutés à la fin.

\begin{decision}[ADR-13 §4, §6 --- un ordre topologique qui n'ordonne rien]
L'ordre effectif de l'analyse est celui des racines (identifiants croissants) et de l'inlining
descendant ; ni \code{analyze_program} ni le driver ne lisent l'ordre topologique. Son unique
consommateur est une ligne de \code{--verbose} (\src{src/driver/mod.rs}{324}{336}) --- c'est la
ligne \og symbol graph: … topo order = […]\fg{} des sorties de ce chapitre. \issue{7} le relevait
déjà. Le reste d'\adr{13} est bien réalisé : clés composites \code{(PathBuf, String)} pour tous
les registres (d'où \code{ComponentKey}), traits \code{FileDiscoverer} et \code{ImportResolver},
entrée par répertoire, \code{FunctionIR} et inlining des utilitaires en position d'instruction,
imports non résolus lus $\Top$ avec une \Info. Le repli \code{get_by_name}, qu'\adr{13} acceptait
comme limite, a été retiré pour les composants par \adr{40}.
\end{decision}

\begin{piege}
La doc de \code{build} affirme un graphe \og over-approximate but never under-approximate\fg.
C'est faux pour un élément construit dans un \code{FnLit} en ligne du rendu (callback de
\code{.map}, render prop) : \code{collect_callees_in_expr} ne traverse pas les \code{FnLit}
(\src{src/engine/symbol_graph.rs}{284}{304}) et \code{CFG::for_each_expr} ne visite que les
expressions de premier niveau. Sur \file{/tmp/ch16/sg1.tsx}
(\code{function Z() { return <ul>{[1].map((i) => <A key={i} />)}</ul>; }}), \code{--verbose}
affiche \og topo order = [Z@sg1.tsx, A@sg1.tsx]\fg{} : aucune arête $Z \to A$, alors que la
détection de racines, qui descend dans les \code{FnLit}, voit bien \code{A} (\og 1 components
analyzed\fg). Sans conséquence sur les résultats, puisque le graphe n'est qu'affiché.
\end{piege}

% ===========================================================================
\section{La dépendance de rendu}
\label{sec:inter-composants-render-deps}

\subsection{Pourquoi une seconde analyse}

\begin{exemple}{Une valeur calculée, une condition, un voisin coûteux}{inter-composants-ctrl}
\begin{lstlisting}[language=TypeScript,caption={\file{/tmp/ch16/ctrl.tsx}},label={lst:inter-composants-ctrl}]
import { useState } from "react";
function Row({ i }: { i: number }) { return <li>row {i}</li>; }
function Heavy() { return <ul>{[1, 2, 3].map((i) => <Row key={i} i={i} />)}</ul>; }
function Hint({ msg }: { msg: string }) { return <p>{msg}</p>; }
function Clear() { return <button>clear</button>; }
export function Search() {
  const [text, setText] = useState("");
  let msg = "";
  if (text.length > 3) msg = "long";
  return (
    <div>
      <input value={text} onChange={(e) => setText(e.target.value)} />
      <Hint msg={msg} />
      {text ? <Clear /> : null}
      <Heavy />
    </div>
  );
}
\end{lstlisting}
\begin{sortie}
$ reactant check ctrl.tsx --rule wasted-subtree-render --trace
  Search  (2 hooks)  ctrl.tsx
    warn   wasted-subtree-render  [hook:0]  (line 12:26)  each `change` event writes state `text` and re-renders `<Heavy>` (at least 2 component renders, including a list) although none of its inputs depends on it. Move `text` and the elements that use it into their own component, or build the unrelated elements higher up and pass them in as `children`
       → `<Heavy>` re-renders (at least 2 component renders, including a list) although none of its props depends on it (line 15:6)
\end{sortie}
\end{exemple}

Chaque frappe re-rend les trois éléments de \code{Search}. \code{<Heavy>} ne dépend de rien :
c'est du gaspillage. \code{<Hint>} reçoit \code{msg}, une chaîne littérale choisie sous une
condition sur \code{text} : sa valeur ne contient pas \code{text}, mais elle en \emph{dépend}.
\code{<Clear>} n'a aucune prop, mais son \emph{existence} dépend de \code{text}. Le domaine des
valeurs ne sait dire aucune de ces deux choses : \code{Stability::Versioned(S)}
(\cref{def:stabilite}) vit sur le seul slot de référence, si bien qu'un booléen,
\code{text.length} ou une allocation \code{{ a: text }} (qui est \code{PerRender}) n'en portent
rien, et \code{MountIndex::Guard} contournait le manque par un balayage syntaxique des
\code{StateVal}.

\begin{decision}[ADR-41 §1 --- une analyse avant séparée, pas un champ de \code{StateValue}]
La dépendance de rendu est une analyse avant, lancée sur un composant \emph{convergé}, à la
demande. Elle n'est \emph{pas} une composante du produit \code{StateValue} : comme champ, elle
changerait toutes les comparaisons de valeurs existantes --- \regle{unnecessary-rerender}
compare \code{arg == init}, la clé du \code{ComponentCache} compare des props --- et la
dépendance est un autre fait que la valeur (un \code{{ a: text }} frais dépend de \code{text}).
Le résumé est \emph{local} au composant (une prop se lit \code{Prop(name)}, quoi que passe le
parent) et se compose sur l'arbre d'éléments à origine prouvée, ce qui garde précis le résumé
d'un composant analysé seulement en phase 2. Cette décision a retiré un contournement
(le balayage syntaxique de \code{MountIndex}, \issue{149}) en corrigeant la cause : c'est le
premier principe du projet appliqué.
\end{decision}

\subsection{Les sources et le treillis \code{Deps}}

\rustsnippet[label={lst:inter-composants-source}]{src/engine/render_deps.rs}{53}{84}{une entrée de rendu, dans le repère d'un composant}

Une source est une entrée du rendu \emph{dans le repère d'un composant} : \code{Slot(0)} du
parent et \code{Slot(0)} de l'enfant sont deux choses différentes, et \code{Prop(name)} est
relatif au composant courant. Seules \code{Context(id)} et \code{Module(name)} sont
\emph{indépendantes du repère} : le même contexte et le même nom de module dans tous les
composants, ce qui leur permet de descendre l'arbre d'éléments sans traduction.

\begin{definition}{Dépendance de rendu}{render-dep}
Soit $\mathit{Src}$ l'ensemble (fini pour un programme) des sources d'un composant. La
\terme{dépendance de rendu} d'une valeur $v$ du rendu est un élément de
$\powerset(\mathit{Src}) \cup \{\Top\}$ ordonné par inclusion, $\Top$ au sommet : un ensemble
\emph{may} des sources dont $v$ peut être calculée, directement (opérande, capture) ou par
contrôle (condition d'une branche qui décide de son affectation). $\Top$ signifie \og peut
dépendre de n'importe quoi\fg. Le join est l'union, $\Top$ absorbant. Une source absente de la
dépendance est une source dont $v$ est \emph{prouvée} indépendante.
\end{definition}

\begin{figure}[htbp]\centering
\begin{tikzpicture}[x=2.3cm,y=1.15cm,treillis,
    n/.style={font=\scriptsize,inner sep=2pt}]
  \node[n] (bot) at (0,0) {$\emptyset$};
  \node[n] (a) at (-1.4,1) {$\{\code{Slot(0)}\}$};
  \node[n] (b) at (0,1) {$\{\code{Prop(v)}\}$};
  \node[n] (c) at (1.4,1) {$\{\code{Module(c)}\}$};
  \node[n] (ab) at (-1.4,2) {$\{\code{Slot(0)},\code{Prop(v)}\}$};
  \node[n] (ac) at (0,2) {$\{\code{Slot(0)},\code{Module(c)}\}$};
  \node[n] (bc) at (1.4,2) {$\{\code{Prop(v)},\code{Module(c)}\}$};
  \node[n] (abc) at (0,3) {$\{\code{Slot(0)},\code{Prop(v)},\code{Module(c)}\}$};
  \node[n] (top) at (0,4) {$\Top$ (\code{top = true})};
  \foreach \u/\v in {bot/a,bot/b,bot/c,a/ab,a/ac,b/ab,b/bc,c/ac,c/bc,ab/abc,ac/abc,bc/abc,abc/top}
    \draw[thin,draw=ra-gray] (\u) -- (\v);
  \node[etiquette,align=left,anchor=west] at (1.9,3.2) {\code{set} : union (may)\\
    \code{gated} $\subseteq$ \code{set} : intersection sur les\\ contributeurs qui portent la source (must)\\
    \code{writes} : union, avec son propre \code{top}};
\end{tikzpicture}
\caption{Le treillis \code{Deps} pour trois sources : parties ordonnées par inclusion, plus un
sommet $\Top$ qui vide \code{set} et \code{gated}. Le treillis est fini, donc la montée
s'arrête sans widening. La composante \code{gated} va dans le sens dual : elle ne peut que
rétrécir quand on joint.}
\label{fig:inter-composants-deps}
\end{figure}

\rustsnippet{src/engine/render_deps.rs}{119}{131}{un ensemble may de sources, et deux compléments}

\code{Deps} porte l'ensemble \code{set}, le drapeau \code{top}, et deux compléments utiles aux
valeurs fonctions. \code{gated} est la partie de \code{set} qu'une fonction n'atteint, quand on
l'appelle, que derrière un test de ses propres arguments (\cref{sec:inter-composants-gated}) ;
\code{writes} est l'ensemble des noms de module qu'appeler la fonction peut écrire
(\code{Writes { top, names }}, \src{src/engine/render_deps.rs}{90}{117}). Leur join cohabite
dans une même fonction :

\rustsnippet[label={lst:inter-composants-union}]{src/engine/render_deps.rs}{153}{170}{le join : may pour \code{set}, must pour \code{gated}}

\code{set} croît par union. Une source reste \code{gated} seulement si \emph{aucun}
contributeur ne l'apporte non gardée : on calcule les sources non gardées des deux côtés, on
réunit les \code{gated}, puis on retire les non gardées. \code{top} absorbe tout, sauf
\code{writes}, qui garde son propre \code{top} (\code{Deps::top()} pose
\code{writes: Writes::top()}, \src{src/engine/render_deps.rs}{141}{147}). Enfin
\code{ungated()} (\src{src/engine/render_deps.rs}{172}{180}) efface \code{gated} en gardant
\code{writes} : appliquée au résultat d'un appel, elle dit que la valeur n'est plus la fonction
décrite (\code{debounce(fn)} rend une fonction qui appelle peut-être \code{fn}), si bien que le
fait may survit et le fait must non.

\begin{piege}
\code{Deps::is_empty()} est \code{!top && set.is_empty()} (\srcline{src/engine/render_deps.rs}{149}) :
il \emph{ignore} \code{writes}. Une fermeture qui n'écrit qu'un binding de module est
\og vide\fg{} pour ce test. Ce n'est pas un défaut --- les questions de rendu passent par
\code{touches}, et les écritures par \code{co_writes} --- mais un appel à \code{is_empty} pour
décider qu'une fonction \og ne fait rien\fg{} serait faux.
\end{piege}

Une question se pose dans un repère sous la forme d'une \code{Relevance { any_prop, sources }}
(\src{src/engine/render_deps.rs}{207}{250}), et la réponse élémentaire est \code{touches}
(\src{src/engine/render_deps.rs}{185}{205}) : une dépendance touche une question si elles
partagent une source, si la dépendance est $\Top$ et la question non vide, ou si l'une porte
l'objet props entier (\code{AllProps}) et l'autre une prop. \code{gated_for(rel)} répond vrai
quand toutes les sources de \code{rel} que la dépendance porte sont gardées
(\cref{chap:regles-rendu-rsc} montre ces deux fonctions).

\subsection{L'analyse avant d'un composant}

\code{render_deps} (\src{src/engine/render_deps.rs}{417}{503}) prend un \code{AnalysisResult}
convergé --- donc le CFG \emph{après} greffe des hooks personnalisés et des utilitaires --- et
l'ensemble \code{written} des noms qu'une fonction quelconque du programme écrit. La valeur
abstraite interne est un \code{DVal { deps, shape }}
(\src{src/engine/render_deps.rs}{333}{379}) : une forme \code{Props} marque l'objet props (un
\code{.field} dessus donne \code{Prop(field)}), une forme \code{Members} garde les sources membre
par membre d'un littéral local (ce dont un retour de hook déstructuré a besoin). L'environnement
d'entrée lie le paramètre props à \code{DVal { {AllProps}, Props }}
(\src{src/engine/render_deps.rs}{447}{454}).

\paragraph{Évaluation.} \code{eval} (\src{src/engine/render_deps.rs}{815}{948}) suit la
\cref{tab:inter-composants-eval}. Deux lignes méritent l'attention : un élément de composant
s'évalue à $\emptyset$ --- ses props sont les entrées de l'enfant, attribuées au site, et
l'élément lui-même n'est pas un usage ---, alors qu'un élément hôte prend l'union de ses props
et de ses enfants.

\begin{table}[htbp]\centering\small
\begin{tabular}{p{42mm}p{100mm}}
\toprule
expression & dépendance \\
\midrule
\code{Lit}, \code{SummaryVal} & $\emptyset$ \\
\code{StateVal(l)}, \code{StateSetter(l)} & \code{{Slot(l)}}, \code{{Setter(l)}} (respectivement) \\
\code{HookMarker(l)} & ref : \code{{Ref(l)}} ; hook opaque : \code{{Hook(l)}} $\cup$ arguments (dé-gardés) ; sinon $\emptyset$ \\
\code{MemoVal(l)}, \code{CallbackVal(l)} & variables libres du corps (fermeture pour un callback) $\cup$ liste de deps ; entrée introuvable : $\Top$ \\
\code{FnLit} & fermeture : captures (gardées si derrière un test des paramètres) $+$ \code{writes} \\
\code{FieldAccess} & sur \code{Props} : \code{{Prop(field)}} ; sur \code{Members} : le membre ; sinon la base \\
\code{IndexAccess}, \code{BinOp} & union des opérandes \\
\code{Call}, \code{New} & callee et arguments, dé-gardés ; un callee \code{use}/\code{useX} lève \code{inline_hook} \\
\code{ObjectLit} & union des champs ; forme \code{Members} si aucun n'est un spread \\
\code{ArrayLit} & union des éléments ; jamais de forme \code{Members} \\
\code{CompApp} & $\emptyset$ \\
\code{NativeElem} & props $\cup$ enfants \\
\bottomrule
\end{tabular}
\caption{L'évaluation de la dépendance de rendu (\src{src/engine/render_deps.rs}{815}{948}).}
\label{tab:inter-composants-eval}
\end{table}

\paragraph{Lecture d'un nom.} Un nom non lié dans le composant --- binding de module, import,
global --- n'est pas forcément une constante :

\rustsnippet{src/engine/render_deps.rs}{603}{625}{un nom non lié : contexte prouvé, binding écrit, ou rien}

Un objet contexte prouvé (\code{ModuleConstInit::Context}) se lit \code{Context(id)} ; un nom que
quelque fonction du programme écrit se lit \code{Module(name)} ; tout le reste est constant d'un
rendu à l'autre et ne dépend de rien. Restreindre \code{Module} aux noms écrits garde les
ensembles petits : sur le dépôt twenty, lire \code{Module} pour \emph{tout} nom libre coûtait
70\,\% (301\,s à 514\,s) ; avec la restriction, l'exécution mesure 366\,s (plan de campagne
\file{docs/campaign/rerender-cascade-plan.md}, §10, \og Corpus effect of module writes\fg).

\paragraph{Transfert et dépendance de contrôle.} \code{run_cfg} itère en ordre RPO jusqu'à
stabilité, au plus \code{MAX_ROUNDS = 64} tours. Au début de chaque bloc, il calcule sa
\emph{condition de chemin} \code{pc} : la dépendance des conditions des branches qui le
contrôlent, évaluées dans l'environnement de sortie du bloc de branche.

\rustsnippet[label={lst:inter-composants-pc}]{src/engine/render_deps.rs}{689}{704}{la condition de chemin d'un bloc, puis le transfert}

Une branche \emph{contrôle} un bloc quand exactement un de ses deux côtés l'atteint : c'est la
différence symétrique des ensembles atteignables depuis \code{then_} et \code{else_}
(\code{controlling_branches}, \src{src/engine/render_deps.rs}{1240}{1276}), la même définition
que \code{MountIndex}. Le transfert (\src{src/engine/render_deps.rs}{777}{813}) ajoute \code{pc} à
toute affectation : \code{env[x] = eval(rhs)} $\cup$ \code{pc}. Dans l'\cref{ex:inter-composants-ctrl},
\code{msg = "long"} est dans le bloc contrôlé par \code{text.length > 3} : \code{msg} prend
\code{{Slot(0)}}, et \code{<Hint>} n'est pas gaspillé. Un \code{MemberWrite obj.k = v} est une
mise à jour \emph{faible} de la racine, qui perd sa forme ; un appel de méthode en instruction
(\code{xs.push(v)}) ajoute ses arguments au récepteur --- sans \code{pc}, et pour toute méthode,
pas seulement les méthodes mutantes. Si la boucle ne converge pas en 64 tours, le drapeau
\code{top} passe à vrai et \emph{tout} le résumé devient $\Top$ : le plafond est une garde, pas
un bouton de précision, et il dégrade vers le haut.

\subsection{La collecte : utiliser ou transmettre}
\label{sec:inter-composants-collecte}

Après le point fixe, une passe de collecte parcourt les blocs dans le même ordre et produit le
résumé :

\rustsnippet[label={lst:inter-composants-renderdeps}]{src/engine/render_deps.rs}{301}{322}{le résumé d'un composant}

C'est le cœur de la dichotomie : \code{genuine} est ce que le composant \emph{utilise},
\code{sites} ce qu'il \emph{transmet}. Entrent dans \code{genuine} : les appels en position
d'instruction, les écritures de membre pendant le rendu et la valeur rendue, chacun avec son
\code{pc} ; tout élément \emph{hôte} (\code{NativeElem}), qui est sortie du composant qui le
construit, même passé en \code{children} à un enfant (\code{<Modal><input value={text}/></Modal>}) ;
les effets (variables libres du corps et liste de deps, évaluées dans l'environnement de sortie
du bloc d'appel, \src{src/engine/render_deps.rs}{466}{497}) ; les arguments d'un hook resté opaque
après expansion. Un \code{Let} ou un \code{Assign} n'est jamais un usage en soi ; une condition de
branche n'en est un qu'à travers le \code{pc} des blocs qu'elle contrôle.

Chaque élément de composant devient un site :

\rustsnippet{src/engine/render_deps.rs}{252}{277}{un élément de composant construit par le rendu}

\rustsnippet{src/engine/render_deps.rs}{961}{979}{la garde d'un site : le chemin, et le type d'élément s'il est une valeur}

La garde d'un site est le \code{pc} de son bloc, plus --- si le nom de type est une valeur du
rendu (\code{const { Modal } = useModal(); <Modal/>}) --- la dépendance de cette valeur : quel
composant se monte en dépend, comme d'une condition (ce correctif a retiré 8 FP sur le dépôt dub
lors de la campagne). Chaque prop explicite garde sa dépendance ; les spreads vont dans
\code{spread} ; \code{provides} est renseigné si le nom, sans \code{.Provider}, désigne un
contexte prouvé non masqué ; \code{parent} est l'indice du site dans lequel celui-ci est imbriqué
comme prop ou enfant. Chaque prop \code{onX} (ou spread) d'un élément hôte devient un
\code{HostHandler} (\src{src/engine/render_deps.rs}{282}{299}), avec la balise, l'attribut
\code{type} littéral et la dépendance de la valeur du handler : c'est là qu'un setter
\og atterrit\fg.

\paragraph{Callbacks synchrones.} Les \code{FnLit} passés à un appel (\code{.map},
\code{.forEach}) tournent pendant le rendu : \code{collect} les analyse par \code{nested}
(\src{src/engine/render_deps.rs}{1079}{1109}), avec des paramètres alimentés par le récepteur et
les autres arguments, \code{in_list} vrai, et une profondeur bornée par \code{MAX_NESTING = 8}
(au-delà, $\Top$). Le corps d'un \code{useMemo} lu pendant le rendu est parcouru de même, sans
\code{in_list}. Un \code{FnLit} qui n'est pas passé à un appel ne tourne pas pendant le rendu :
seules ses captures comptent, comme dépendance de la valeur fonction.

\begin{exemple}{Une écriture dans un \code{forEach} est perdue}{inter-composants-foreach}
\begin{lstlisting}[language=TypeScript,caption={\file{/tmp/ch16/foreach.tsx} (extrait)},label={lst:inter-composants-foreach}]
export function App() {
  const [sel, setSel] = useState(0);
  let label = "";
  [1, 2, 3].forEach((i) => { if (i === sel) label = String(i); });
  return (
    <div>
      <input value={sel} onChange={(e) => setSel(Number(e.target.value))} />
      <Child label={label} />
      <Heavy />
    </div>
  );
}
\end{lstlisting}
\begin{sortie}
$ reactant check foreach.tsx --rule wasted-subtree-render
  App  (2 hooks)  foreach.tsx
    warn   wasted-subtree-render  [hook:0]  (line 10:25)  each `change` event writes state `sel` and re-renders `<Child>` and `<Heavy>` (2 component renders) although none of their inputs depends on it. Move `sel` and the elements that use it into their own component, or build the unrelated elements higher up and pass them in as `children`
\end{sortie}
L'affirmation est fausse pour \code{<Child>} : \code{label} est calculé à partir de \code{sel}.
Deux faits du code se combinent. \code{nested} lance \code{run_cfg(body, …, false)} et
\emph{jette} l'environnement de sortie du callback : une écriture du callback dans un binding
local du rendu ne remonte pas. Et \code{compute_free_vars} retire \code{label}, cible d'un
\code{Assign}, des variables libres du callback (\src{src/ir/free_vars.rs}{96}{106}). \code{label}
garde donc la dépendance de \code{""}, $\emptyset$. C'est une source manquante, donc un faux
positif d'une règle d'absence --- la direction qu'\adr{41} §2 veut exclure ; il n'était ni
documenté dans \file{docs/limitations.md} ni suivi par une issue au snapshot.
\end{exemple}

\subsection{Handlers gardés : le fait \code{gated}}
\label{sec:inter-composants-gated}

Une frappe dans un champ texte écrit à chaque touche ; un \code{onKeyDown} qui n'écrit que sur
\code{Enter} écrit une fois par validation. La différence n'est pas une question de preuve mais
de \emph{fréquence}, et la règle \regle{wasted-subtree-render} ne signale par défaut que les
déclencheurs continus (option \code{continuousOnly}, \adr{41}, \cref{chap:regles-rendu-rsc}).
\code{param_gated_vars} (\src{src/engine/render_deps.rs}{1176}{1238}) calcule, pour une fonction,
les variables libres qu'elle n'atteint que derrière un test de ses paramètres : il propage les
variables \emph{teintées} par les paramètres jusqu'à un point fixe, et les blocs contrôlés par
une condition qui en lit une ; une variable libre est gardée si elle n'apparaît que dans de tels
blocs. Sur la fixture \file{tests/fixtures/wasted_subtree_render/keyed.tsx}, le handler de
\code{Submit} (\code{if (e.key !== "Enter") return; setQ(...)}) a la dépendance
\code{{Setter(0)}} avec \code{gated = {Setter(0)}} (résumé relevé par la sonde du dossier 08) ;
\code{KeyLog} (\code{onKeyDown={(e) => setLast(e.key)}}) ne garde rien. Par défaut, seul
\code{KeyLog} est signalé ; avec \code{continuousOnly=false}, \code{Escape}, \code{Handed},
\code{KeyLog} et \code{Submit} le sont (observé). Le fait est \emph{must} : un seul contributeur
non gardé l'annule (\cref{lst:inter-composants-union}), et il ne peut que classer un
déclencheur comme discret, jamais changer une preuve.

\subsection{Contextes}
\label{sec:inter-composants-contextes}

\tsxsnippet[label={lst:inter-composants-context}]{tests/fixtures/state_lifted_too_high/context.tsx}{6}{21}{un état qui atteint son consommateur par contexte}

Dans le résumé d'\code{App} (sonde du dossier 08), le provider est un site
\code{provides = Some(Query)} dont la prop \code{value} a la dépendance
\code{{Slot(0), Setter(0)}} ; \code{<Header>} et \code{<Page>} ont \code{parent = Some(0)} ;
\code{children} vaut $\emptyset$, puisqu'un élément de composant s'évalue à $\emptyset$. Dans
\code{Search}, \code{useContext(Query)} est un hook opaque : sa dépendance est \code{Hook(0)}
$\cup$ celle de son argument, \code{Context(Query)}. Le drapeau \code{any_context} reste faux,
parce que \code{Query} est prouvé et écrit dans le composant ; il passe à vrai pour un
\code{useContext} d'un objet non prouvé ou atteint par un hook inliné, et pour tout hook de code
utilisateur que l'analyse n'a pas pu inliner (\code{reads_unnamed_context},
\src{src/engine/render_deps.rs}{645}{670}). Un hook de paquet, lui, ne peut atteindre un contexte
des modules utilisateur que si on le lui passe, ce que ses arguments montrent.

\begin{sortie}
$ reactant check tests/fixtures/state_lifted_too_high/context.tsx --rule state-lifted-too-high --trace
  App  (1 hooks)  tests/fixtures/state_lifted_too_high/context.tsx
    warn   state-lifted-too-high  [hook:0]  (line 14:8)  state `q` is only used inside `<Search>`, 2 levels below `App`. Every write re-renders `App`, `Page` and 2 other components they render although none of them uses it; move the state into `Search`, with the `Query` provider that hands it on
       → `App` renders `<Page>` inside the `Query` provider without using it itself (line 18:6)
       → `Page` renders `<Search>` inside the `Query` provider without using it itself (line 12:41)
\end{sortie}

Côté domaine des valeurs, \code{Query.Provider} est un composant \emph{inconnu} du registre : il
compte dans \code{unknown_component_refs}, d'où les \Info{} \regle{analysis-limit} sur les
\code{X.Provider}, et c'est le cas courant du trou de havoc de l'\cref{ex:inter-composants-hv}.
Dans \code{Guarded}, plus bas dans la même fixture, un \code{<Tooltip/>} non résolu sous le
provider compte comme consommateur possible : pas de finding.

\subsection{Bindings de module}

Un handler peut changer autre chose qu'un état : \code{cache.hits += 1}, \code{counter++},
\code{seen.set(k, v)}. Si un autre composant lit ce binding, il n'est pas gaspillé.
\code{written_names} (\src{src/engine/render_deps.rs}{406}{412}) rassemble les racines écrites du
rendu et de chaque corps de hook ; \code{written_roots} (\src{src/engine/render_deps.rs}{1125}{1171})
retient la cible d'un \code{Assign} (l'IR épelle ainsi l'écriture d'un binding extérieur), la
racine d'un \code{MemberWrite}, le récepteur d'une méthode mutante (la liste d'\adr{28} :
\code{push}, \code{pop}, \dots, \code{set}, et \code{Object.assign}), fermetures imbriquées
comprises, moins ce que le corps lie. L'union sur tous les composants est l'ensemble
\code{written} de la lecture d'un nom. Puis \code{writes_of} (\src{src/engine/render_deps.rs}{580}{601})
traduit les racines écrites d'un corps en noms de module \emph{à travers les valeurs qu'elles
peuvent aliaser} (\code{const c = cache; c.x = 1} écrit \code{cache}) ; une racine de valeur
$\Top$ donne \code{Writes::top()}.

\begin{sortie}
$ reactant check tests/fixtures/wasted_subtree_render/module_write.tsx --rule wasted-subtree-render --rule-option wasted-subtree-render:continuousOnly=false
  Control  (2 hooks)  …  each `change` event writes state `text` and re-renders `<Hits>`, `<Count>`, `<Seen>`, `<Stats>` and `<Chart>` (6 component renders) …
  Debounced  (2 hooks)  …  each `change` event writes state `text` and re-renders `<Chart>` (2 component renders) …
  Handed  (1 hooks)  …  each `change` event in `<Field>` writes state `text` and re-renders `<Chart>` (2 component renders) …
  Typing  (2 hooks)  …  each `change` event writes state `text` and re-renders `<Chart>` (2 component renders) …
\end{sortie}

(lignes élaguées aux en-têtes et au début des messages.) \code{Control} n'écrit que son état :
les cinq voisins sont gaspillés. Le handler de \code{Typing} écrit \code{cache}, \code{counter}
et \code{seen} à côté de \code{setText} : \code{Hits}, \code{Count}, \code{Seen} et \code{Stats}
(qui lit \code{cache} à travers l'utilitaire \code{describe()} inliné) cessent d'être gaspillés,
seul \code{Chart} reste. Dans \code{Handed}, le setter descend vers \code{Field}, dont la
fermeture écrit \code{last} : \code{writes} s'accumule le long du trajet. Dans \code{Debounced},
\code{debounce(...)} est opaque, mais \code{ungated()} garde \code{writes}.

\begin{remarque}
Les noms de module s'apparient \emph{par orthographe}, à travers les fichiers : deux
\code{cache} sans rapport n'en font qu'un, ce qui donne moins de findings, jamais un faux. Une
hypothèse reste, énoncée par \adr{41} §2 : un appel que l'analyse ne voit pas ne lit ni n'écrit
de binding de module (\issue{51}, \issue{52}).
\end{remarque}

% ===========================================================================
\section{Composer sur l'arbre d'éléments}
\label{sec:inter-composants-arbre}

\subsection{L'index de rendu et la résolution d'un site}

Les résumés sont locaux ; les questions sont globales (\og où cet état est-il utilisé ?\fg). La
couche règles les compose dans \code{RenderIndex} (\file{src/rules/helpers/render_tree.rs}),
construit une fois par programme : \code{build} (\src{src/rules/helpers/render_tree.rs}{117}{137})
calcule \code{written}, puis un résumé par composant, puis le nombre de sites qui nomment chaque
composant. Passer d'un site à un composant se fait par \code{resolve}
(\src{src/rules/helpers/render_tree.rs}{767}{777}), qui n'accepte que l'origine prouvée
(\code{table.id_of(origin)}).

\begin{piege}
Il y a donc deux notions de \og résolu\fg. Le moteur (\code{resolve_child}) accepte l'origine,
sinon un nom unique ; l'arbre d'éléments n'accepte que l'origine. Un site sans origine dont le
nom est unique est inliné par le moteur mais opaque pour les cascades. Chacun est prudent dans
sa polarité : l'inliner qui trouverait un enfant par son nom analyse au pire un corps de trop,
alors que la règle de cascade qui se tromperait d'enfant \emph{cacherait} les usages du vrai.
\end{piege}

\subsection{\code{uses} : descendre en traduisant la question}

\code{uses(comp, rel)} (\src{src/rules/helpers/render_tree.rs}{516}{592}) construit un arbre
d'usages. Le nœud est \emph{utilisateur} si le résumé \code{uses(rel)} (\code{genuine} touche la
question, ou \code{any_context} et la question nomme un contexte). Pour chaque site : une garde
touchée fait du nœud un utilisateur ; un provider est sauté, sa valeur étant suivie par contexte ;
sinon on calcule les props qui portent la question (\code{forwarded}, qui laisse passer un spread
de l'objet props sous les mêmes noms et traite tout autre spread comme \og n'importe quelle
prop\fg), les contextes qui la portent (\code{contexts_at}, qui remonte les \code{parent}
jusqu'aux providers dont la \code{value} la touche) et les modules. Un site dans une liste, non
résolu, trop profond (\code{MAX_DEPTH = 64}) ou déjà visité fait du nœud un utilisateur. Sinon on
descend avec la question \emph{traduite dans le repère de l'enfant} : \code{Prop(p)} pour chaque
prop qui la porte, \code{any_prop} si un spread peut renommer, contextes et modules tels quels.

\subsection{\code{home_of} : le plus petit sous-arbre}

\code{home_of(owner, label)} (\src{src/rules/helpers/render_tree.rs}{195}{257}) rend \code{None}
si un écrivain du slot peut écrire n'importe quoi (\code{co_writes.top}), si personne ne lit le
slot ou si personne ne l'écrit. Sinon il construit l'arbre d'usages de la question
\code{{Slot(l), Setter(l)}} $\cup$ \code{co_writes} et descend tant qu'exactement un enfant a
des usages :

\rustsnippet{src/rules/helpers/render_tree.rs}{226}{244}{descendre tant qu'un seul enfant utilise}

Le chemin parcouru est la liste des composants qui re-rendent pour rien ; \code{siblings} compte
les autres sites (hors providers) que construisent les composants du chemin.
\code{co_writes} (\src{src/rules/helpers/render_tree.rs}{154}{178}) réunit les \code{writes} des
\emph{landings} du setter et les \code{effect_writes} des effets qui écrivent le slot ;
\code{landings} (\src{src/rules/helpers/render_tree.rs}{379}{512}) suit la capacité
\code{Setter(l)} de prop en prop jusqu'aux handlers hôtes dont la dépendance la touche --- une
prop à la fois, car \og a prop that lands below must not hide one that does not\fg{} ---, et une
prop \code{onX} qui n'atterrit nulle part plus bas atterrit sur l'élément lui-même.

\begin{figure}[htbp]\centering
\begin{tikzpicture}[node distance=9mm and 12mm,
    comp/.style={bloc,minimum width=24mm},
    chemin/.style={comp,draw=ra-amber,very thick},
    maison/.style={comp,draw=ra-blue,fill=ra-blue!12,very thick}]
  \node[chemin] (app) {\code{App}\\{\scriptsize possède \code{Slot(0)}, \code{Setter(0)}}\\{\scriptsize \code{genuine} $=\emptyset$}};
  \node[chemin,below=of app] (lay) {\code{Layout}\\{\scriptsize \code{genuine} $=\emptyset$}};
  \node[chemin,below left=of lay] (side) {\code{Sidebar}\\{\scriptsize \code{genuine} $=\emptyset$}};
  \node[comp,below right=of lay,draw=ra-gray] (cont) {\code{Content}\\{\scriptsize aucune prop : frère}};
  \node[maison,below=of side] (field) {\code{Field}\\{\scriptsize \code{genuine} $=$ \{\code{Prop(value)},}\\{\scriptsize \code{Prop(onChange)}\}}};
  \draw[fleche] (app) -- node[etiquette,right,align=left] {\{\code{Prop(value)}, \code{Prop(onChange)}\}} (lay);
  \draw[fleche] (lay) -- node[etiquette,left,align=right] {\{\code{Prop(value)},\\ \code{Prop(onChange)}\}} (side);
  \draw[fleche,draw=ra-gray] (lay) -- node[etiquette,right] {$\emptyset$} (cont);
  \draw[fleche] (side) -- node[etiquette,left,align=right] {\{\code{Prop(value)},\\ \code{Prop(onChange)}\}} (field);
  \node[etiquette,right=4mm of field,align=left] {\emph{home} : le handler de l'\code{<input>}\\ lit \code{Prop(onChange)}, sa \code{value}\\ lit \code{Prop(value)}};
\end{tikzpicture}
\caption{L'arbre d'éléments de \code{drill.tsx} et la question \code{{Slot(0), Setter(0)}}
d'\code{App}, traduite à chaque saut dans le repère de l'enfant. Les trois composants du chemin
(bord orangé) ne font que transmettre ; \code{Field} (bleu) est le \emph{home} ;
\code{Content} est un frère qui re-rend aussi (\code{siblings = 1}).}
\label{fig:inter-composants-drill}
\end{figure}

\begin{exemple}{Le home de \code{value} dans \code{drill.tsx}}{inter-composants-home}
Les résumés (sonde du dossier 08) : \code{App}, \code{Layout} et \code{Sidebar} ont
\code{genuine} $= \emptyset$, et chacun transmet \code{value} et \code{onChange} à son unique
enfant utilisateur ; \code{Field} a \code{genuine} $=$ \code{{Prop(onChange), Prop(value)}}
et un handler \code{change} sur \code{<input>} de dépendance \code{{Prop(onChange)}}.
\code{home_of(App, 0)} descend \code{App} $\to$ \code{Layout} $\to$ \code{Sidebar} $\to$
\code{Field} (\cref{fig:inter-composants-drill}), avec un frère, \code{Content} :
\begin{sortie}
$ reactant check tests/fixtures/state_lifted_too_high/drill.tsx --rule state-lifted-too-high --trace
  App  (1 hooks)  tests/fixtures/state_lifted_too_high/drill.tsx
    warn   state-lifted-too-high  [hook:0]  (line 17:8)  state `value` is only used inside `<Field>`, 3 levels below `App`. Every write re-renders `App`, `Layout`, `Sidebar` and 1 other component they render only to pass it down; move the state into `Field`
       → `App` passes it to `<Layout>` as `value`, `onChange` without using it itself (line 18:9)
       → `Layout` passes it to `<Sidebar>` as `value`, `onChange` without using it itself (line 14:15)
       → `Sidebar` passes it to `<Field>` as `value`, `onChange` without using it itself (line 11:32)
\end{sortie}
La règle elle-même, ses seuils et ses options sont au \cref{chap:regles-etat}.
\end{exemple}

\code{wasted_siblings(owner, rel)} (\src{src/rules/helpers/render_tree.rs}{266}{326}) répond à la
question complémentaire, celle de \regle{wasted-subtree-render} (\cref{chap:regles-rendu-rsc}) :
les sites qu'aucune garde, prop ou spread (hors provider), ni aucun site englobant, ne relie à la
question ; seuls les composants résolus sont candidats, un enfant qui utilise un contexte porté
ou un module écrit est écarté, et le coût est une borne inférieure (une liste compte une fois).

\begin{remarque}
Tout le fichier \file{render_tree.rs} obéit à la règle que son en-tête énonce
(\src{src/rules/helpers/render_tree.rs}{11}{13}) : l'absence d'usage y est une preuve, donc tout
inconnu est un usage --- élément non résolu, récursion, profondeur, résumé $\Top$, contexte non
nommé, écrivain de racine $\Top$ (\adr{41} §2). Le \cref{chap:regles-rendu-rsc} en fait un
invariant.
\end{remarque}

% ===========================================================================
\section{Soundness : ce qui est garanti, ce qui fuit}
\label{sec:inter-composants-soundness}

Les deux mécanismes du chapitre n'ont pas la même polarité. L'inlining alimente des règles qui
prouvent une \emph{présence} (un setter appelé, une boucle) : il doit sur-approximer l'état. La
dépendance de rendu alimente des règles qui prouvent une \emph{absence} (\og aucune entrée ne
dépend de l'état\fg) : pour elles, une source manquante est un faux positif, et \og sound\fg{}
veut dire \og tout inconnu est un usage\fg. La \cref{tab:inter-composants-soundness} récapitule
les garanties et les fuites vérifiées au snapshot.

\begin{table}[htbp]\centering\small
\begin{tabular}{p{40mm}p{48mm}p{52mm}}
\toprule
mécanisme & garantie & fuite vérifiée \\
\midrule
résolution (\code{resolve_child}) & jamais de devinette : ambigu $=$ inanalysable & repli sur le nom écrit après un barrel (précision seulement) \\
enfant inanalysable & havoc des setters échappés & setter champ d'un littéral inline (\cref{ex:inter-composants-hv}) \\
récursion & coupure déclarée, assurances suspendues & aucun havoc (\cref{ex:inter-composants-rec}) \\
store partagé & monotone, importé par join & --- \\
cache & saut d'analyse sound pour le store & \code{results[child]} écrasé : FN non déclaré (\cref{ex:inter-composants-overwrite}) \\
élément argument d'un hook opaque & racine en trop, \Info{} sur le hook & ni inliné ni havoqué : \code{App: 0 iteration(s)} sur \file{hookarg.tsx} \\
phase 2 & \code{phase1_reached}, \code{complete_ancestry} & plafonds non enregistrés hors \code{InterCtx} \\
dépendance de rendu & plafonds $\Rightarrow \Top$ ; inconnu $=$ usage & écriture d'un callback synchrone perdue (\cref{ex:inter-composants-foreach}) \\
bindings de module & \code{Module(name)}, \code{writes} & appel opaque supposé muet (\issue{51}, \issue{52}) \\
\bottomrule
\end{tabular}
\caption{Garanties et fuites du sous-système inter-composants.}
\label{tab:inter-composants-soundness}
\end{table}

La ligne \og élément argument d'un hook opaque\fg{} demande un mot. Sur
\file{/tmp/ch16/hookarg.tsx} (\code{const m = useModal(<Dialog onClose={setS} />)}, avec un
\code{Dialog} qui appelle \code{onClose(1)} pendant son rendu), \code{--verbose} rapporte
\og cache hits: 0, misses: 0\fg{} et \og App: 0 iteration(s)\fg : l'élément argument n'est jamais
évalué par le domaine des valeurs, donc ni inliné, ni havoqué, et \code{s} reste à sa valeur
initiale, alors qu'un \code{useModal} qui rend l'élément ferait appeler \code{setS}. Le faux
négatif est déclaré par l'\Info{} du hook inconnu (\og hook \code{`useModal`} was not found in the
registry … (FN possible)\fg), qui suspend les assurances d'\code{App}. Vérifié au
\snapshotCommit{} au-delà des seuls compteurs : \code{make_hook_entry}
(\src{src/lowering/hook_extractor.rs}{641}{660}) retire les arguments d'un hook Custom du CFG et
les range dans \code{HookEntry::Custom::args}, sans laisser de \code{Expr::Call} qui les
contiendrait encore ; \code{expand_custom_hooks} (\src{src/engine/fixpoint.rs}{889}{992}) ne les
lit que pour une substitution syntaxique (\code{param_subst}, seuls les inits \code{State} du
hook expansé) ; \code{custom_arg_returns} (\src{src/engine/fixpoint.rs}{240}{275}) n'exécute
qu'un argument \code{FnLit} ou une variable dont la liaison est certifiée fonction
(\code{certified_fn_binding}) --- un \code{Expr::CompApp} ne correspond à aucun des deux et
retombe dans le \code{continue}. Aucun autre point du moteur n'évalue les arguments d'un hook
Custom : le \code{Dialog} de l'exemple n'est donc bien ni inliné ni havoqué.

Enfin, trois commentaires de code sont en retard sur lui : \code{components_analyzed}
(\og including re-analyses\fg), \code{recursive_component_refs} et
\code{recursive_components} (\og $\Top$\fg{} au lieu de \code{Stable}), et la doc du
\code{ComponentCache} (\og sound over-approximation\fg{} pour une entrée dont le résultat est
celui de la dernière insertion).

% ===========================================================================
\begin{resume}
\begin{itemize}
  \item Un analyseur intra lit les props comme $\Top$ : sound, mais aveugle aux défauts qui
        traversent une frontière (un setter du parent appelé par l'enfant). L'analyse
        inter-composants (\adr{12}) inline chaque enfant au point où le parent construit
        l'élément, avec les props abstraites du parent, et fait remonter ses écritures.
  \item \textbf{Identité} (\adr{40}) : clé \code{(file, name)}, \code{ComponentId} interné par
        le registre dans une \code{ComponentTable} portée par le résultat ; le nom d'affichage
        (\code{Name@file} en cas de collision) n'est qu'un rendu, jamais une clé
        (\code{RuleCtx::component()} contre \code{component_name()}). Un élément porte
        l'origine que son fichier a prouvée ; \code{resolve_child} répond \code{Resolved},
        \code{Unknown} ou \code{Ambiguous}, et refuse de deviner.
  \item \textbf{Racines} : heuristique (ce qu'aucun \code{CompApp} ne référence, une référence
        non résolue comptant pour tous les homonymes), \code{--all-roots}, \code{--entry}
        (un nom qui ne sélectionne rien est une erreur d'usage).
  \item \textbf{Inlining} : \code{eval_comp_app} résout, coupe la récursion (sans havoc),
        évalue les props dans le tas, consulte le cache, analyse l'enfant sous un
        \code{InterCtx} fils, range son résultat. Un élément vaut toujours \code{Stable} ; un
        \code{return} direct est évalué deux fois par passe.
  \item \textbf{Flux ascendant} : un setter porte son propriétaire ; l'appeler écrit le
        \code{SharedStateStore}, monotone, que la boucle du parent importe à chaque test de
        convergence ; pas de point fixe de programme. Un setter qui s'échappe vers un enfant
        inanalysable est havoqué (trou : setter champ d'un littéral inline).
  \item \textbf{Cache} : égalité de treillis, cinq entrées, éviction par join. Ce n'est qu'un
        \og déjà vu\fg{} : \code{results[child]} est le dernier raté, jamais un join, d'où un FN
        dépendant de l'ordre des racines.
  \item \textbf{Deux phases} : racines sous \code{InterCtx}, puis balayage intra ;
        \code{phase1_reached} photographié entre les deux ; hors de lui, une ascendance vide est
        inconnue (\code{complete_ancestry}).
  \item \textbf{Graphe de symboles} (\adr{13}) : construit, trié, affiché en \code{--verbose},
        n'ordonne rien.
  \item \textbf{Dépendance de rendu} (\adr{41}, \cref{def:render-dep}) : analyse avant séparée
        sur le composant convergé ; \code{Deps} $= \powerset(\mathit{Src}) \cup \{\Top\}$, may,
        avec \code{gated} (must) et \code{writes} ; dépendance de contrôle par les branches qui
        contrôlent un bloc ; \code{genuine} (utilisé) contre \code{sites} (transmis) ; handlers,
        contextes (\code{Context(id)}, \code{any_context}), bindings de module
        (\code{Module(name)}). Composition sur l'arbre à origine prouvée : \code{uses},
        \code{home_of}, \code{landings}, \code{wasted_siblings} ; tout inconnu est un usage.
  \item Fichiers : \file{component_registry.rs}, \file{root_detector.rs},
        \file{component_cache.rs}, \file{symbol_graph.rs}, \file{analysis_result.rs},
        \file{program_result.rs}, \file{render_deps.rs} ; \file{component_id.rs},
        \file{context.rs}, \file{state_value.rs}, \file{fixpoint.rs} ;
        \file{rules/helpers/render_tree.rs}.
\end{itemize}
La partie suivante passe aux règles : le \cref{chap:api-regles} montre comment elles lisent ces
résultats par des requêtes typées et un sceau de sévérité ; \regle{state-lifted-too-high} et
\regle{wasted-subtree-render} (\cref{chap:regles-etat,chap:regles-rendu-rsc}) sont les deux
consommateurs de la dépendance de rendu.
\end{resume}

\section*{Exercices}
\addcontentsline{toc}{section}{Exercices}
\label{sec:inter-composants-exercices}

\begin{exercice}{Compter les passages dans le cache}{inter-composants-drill-cache}
Sur \file{drill.tsx} (\cref{lst:inter-composants-drill}), \code{--verbose} affiche
\og cache hits: 3, misses: 7\fg{} et \og App: 1 iteration(s)\fg. Retrouver ces nombres à partir
des trois faits du \cref{sec:inter-composants-evaluations}.
\end{exercice}
\begin{solution}
\code{App} fait deux passes (une itération, due au setter appelé par \code{Field} via le store
partagé). Son \code{return <Layout …/>} est direct : deux évaluations par passe. Passe 1 : raté
(\code{Layout} analysé, donc \code{Sidebar}, \code{Field} et \code{Content}), puis succès. Passe 2 :
\code{value} a changé, raté (\code{Layout} ré-analysé, et avec lui \code{Sidebar} et
\code{Field}, dont les props ont changé), puis succès. \code{Content} (props vides) rate la
première fois et réussit à la seconde analyse de \code{Layout}. Ratés : 2 (\code{Layout}) + 2
(\code{Sidebar}) + 2 (\code{Field}) + 1 (\code{Content}) = 7 ; succès : 2 (\code{App}) + 1
(\code{Content}) = 3.
\end{solution}

\begin{exercice}{Pourquoi tous les homonymes}{inter-composants-homonymes}
Dans l'heuristique des racines, une référence sans origine marque comme référencés
\emph{tous} les composants de ce nom. Que se passerait-il si elle n'en marquait aucun ? si elle
n'en marquait qu'un, le premier par tri ?
\end{exercice}
\begin{solution}
N'en marquer aucun ferait de chaque homonyme une racine analysée avec des props $\Top$ ; si l'un
d'eux est aussi inliné par un parent qui l'importe, l'ordre des identifiants peut faire gagner
le résultat imprécis (écrasement de \code{results}). N'en marquer qu'un reproduirait la devinette
qu'\adr{40} a supprimée : la référence ne désigne pas forcément celui-là. Marquer tous les
homonymes ne coûte que des racines en moins, jamais une conclusion.
\end{solution}

\begin{exercice}{Corriger l'écrasement à la racine}{inter-composants-overwrite}
Reproduire l'\cref{ex:inter-composants-overwrite}, puis proposer une correction de
\code{results[child]} qui se justifie en un paragraphe (\og règle du paragraphe unique\fg).
\end{exercice}
\begin{solution}
Deux pistes. Indexer les résultats par (identité, props abstraites) --- la clé du cache --- et
faire tourner les règles intra-composant sur chaque contexte d'un composant : c'est la
sensibilité au contexte qu'\adr{12} §2 revendique, poussée jusqu'aux règles ; le coût est un
multiplicateur borné par la taille du cache. Ou joindre les résultats d'un même composant, ce
qui exige un join d'\code{AnalysisResult} qui n'existe pas (CFG, relations et témoins ne se
joignent pas comme des états, \cref{sec:inter-composants-cache}). La première se justifie en une
phrase (\og une règle doit voir chaque contexte d'appel analysé\fg) ; la seconde non. Aucune n'est
implémentée au snapshot.
\end{solution}

\begin{exercice}{Perdre le fait \code{gated}}{inter-composants-gated}
Écrire un handler \code{onKeyDown} qui appelle \code{setQ} derrière \code{if (e.key === "Enter")}
\emph{et} ailleurs sans test. Que devient \code{gated} ? Relier la réponse à
\cref{lst:inter-composants-union}.
\end{exercice}
\begin{solution}
Par exemple
\code{(e) => { if (e.key === "Enter") setQ(e.currentTarget.value); log(setQ); }}.
\code{param_gated_vars} trouve \code{setQ} dans un bloc contrôlé par un test de \code{e}
\emph{et} dans un bloc non contrôlé : \code{setQ} est dans \code{open}, donc pas gardé. Plus
généralement, dans \code{union_with}, une source apportée non gardée par un seul contributeur
est retirée de \code{gated} : le fait est must. Le déclencheur redevient continu. Observé sur
\file{/tmp/ch16/gated.tsx}, avec un \code{<Heavy/>} voisin : sans option,
\regle{wasted-subtree-render} ne signale que le composant \code{Mixed} (\og each \code{`keydown`}
event writes state \code{`q`} and re-renders \code{`<Heavy>`}\fg), pas \code{Keyed}, dont le seul
appel est gardé.
\end{solution}

\begin{exercice}{Repères}{inter-composants-frame}
Expliquer pourquoi \code{Module(name)} et \code{Context(id)} descendent l'arbre d'éléments sans
traduction, alors que \code{Slot(l)} et \code{Prop(p)} doivent être traduits à chaque saut.
\end{exercice}
\begin{solution}
\code{Slot(l)} désigne le slot $l$ \emph{du composant courant} : le même label dans l'enfant est un
autre slot. \code{Prop(p)} désigne la prop $p$ \emph{reçue} : la prop \code{value} du parent
devient, chez l'enfant, la prop sous laquelle il la transmet. Un nom de module ou un contexte
prouvé désigne le même objet dans tous les composants (appariement par orthographe pour le
premier, \code{ContextId} canonique pour le second) : \code{uses} les recopie tels quels dans la
question traduite.
\end{solution}

\begin{exercice}{Un havoc complet}{inter-composants-havoc}
Sur l'\cref{ex:inter-composants-hv}, où faut-il corriger pour que \code{B} et \code{D} soient
havoqués, et pourquoi pas dans une règle ?
\end{exercice}
\begin{solution}
Dans \code{collect_escaping_setters}, au niveau où il descend dans les champs d'un
\code{ObjectLit} ou les éléments d'un \code{ArrayLit} : un champ dont la \emph{valeur} est un
setter (\code{as_setter()}) doit être poussé, comme l'est le setter nu au niveau de la prop.
C'est le mécanisme partagé de soundness du domaine des valeurs : toutes les règles lisent l'état
qu'il calcule, et un correctif par règle serait le contournement que le premier principe du
projet interdit.
\end{solution}

\begin{exercice}{Ce qui manque à \code{nested}}{inter-composants-foreach}
Sur l'\cref{ex:inter-composants-foreach}, quel fait faudrait-il remonter de l'analyse du
callback pour que \code{<Child>} ne soit pas signalé ? Pourquoi la copie des seules variables
libres aggrave-t-elle le problème ?
\end{exercice}
\begin{solution}
Il faudrait joindre, dans l'environnement du rendu, la dépendance des bindings \emph{extérieurs}
que le callback affecte (ici \code{label}, avec le \code{pc} du test \code{i === sel}). Or
\code{nested} jette l'environnement de sortie, et \code{compute_free_vars} retire les cibles
d'\code{Assign} des variables libres : \code{label} n'est même pas copié dans l'environnement du
callback, et rien ne dit qu'il est écrit. Les racines écrites existent pourtant déjà
(\code{written_roots}) : c'est le fait à réutiliser.
\end{solution}

\begin{exercice}{Le bottom qui n'est pas neutre}{inter-composants-exit-env}
Pourquoi \code{AnalysisResult::exit_env} joint-il les environnements de sortie avec
\code{reduce} et non avec \code{fold(AbstractEnv::bottom(), join)} ?
\end{exercice}
\begin{solution}
Dans le join d'environnements, une clé présente d'un seul côté prend $\Top$
(\cref{chap:transfert-stores}). Partir de l'environnement vide ferait donc passer à $\Top$
toute variable de l'autre environnement : $\Bot$ n'est pas l'élément neutre du join. \code{reduce}
part du premier environnement réel, et seul le cas \og aucun bloc \code{Return}\fg{} rend
l'environnement vide.
\end{solution}
